% ---------------------------------------------------------------------------
% Chapitre 13 — Le point fixe d'un composant
% Sources : notes/06-engine-fixpoint.md (intégralement, hors §4.4, §4.9,
% §4.10 traités au chapitre 12), notes/16-react-semantics.md (§4.1, §5.1),
% src/engine/fixpoint.rs (analyze_component_impl et ce qu'il appelle),
% src/engine/triggers.rs, src/engine/hook_registry.rs,
% src/domains/stores/{state_store,memo_store,shared_state_store}.rs,
% src/engine/analysis_result.rs, src/rules/impls/infinite_loop.rs ;
% ADR-001, 004, 005, 009, 012, 014, 017, 042.
% Sorties CLI observées au commit e67b10a avec target/debug/reactant sur
% /tmp/ch13/*.tsx. Les valeurs internes (stores, widen_trace,
% effect_setter_writes, effect_triggers, SharedStateStore) ont été relevées
% avec la sonde /tmp/eng/probe (crate hors dépôt liant `reactant` par chemin),
% qui appelle analyze_component et analyze_program et imprime les champs de
% l'AnalysisResult ; ses tableaux de tours viennent d'une réplique de la
% boucle externe dont le résultat final coïncide avec celui du moteur sur
% tous les exemples cités (sauf mention contraire).
% ---------------------------------------------------------------------------
\chapter{Le point fixe d'un composant}
\label{chap:point-fixe-composant}

\begin{objectifs}
  \item Comprendre ce que calcule \code{analyze_component_impl} : non pas une simulation de React, mais un sur-ensemble des états qu'un composant peut atteindre, sous la forme d'un post-point fixe du \code{StateStore}.
  \item Connaître la préparation d'un composant (constantes de module, inlining des utilitaires et des hooks personnalisés, seuils de widening, amorçage des \code{useState}) et ce qu'elle laisse dans les CFG.
  \item Savoir dérouler un tour de la boucle externe --- rendu, mémos, effets, handlers, join, test --- et expliquer chaque choix par un fait de React : batching, fermetures, deps, événements.
  \item Maîtriser la convergence : le test $N \lleq S$, le compteur global, le widening à seuils, \code{widen_trace}, le garde-fou à cent tours, et pourquoi tout cela termine.
  \item Connaître les deux passes post-convergence (rafraîchissement du rendu, \code{effect_setter_writes}) et la tranche où naissent les relations.
  \item Situer la boucle par rapport à React-tRace, en donner l'argument de soundness, et connaître, un par un, les trous observés au commit \snapshotCommit.
\end{objectifs}

\prerequis{\cref{chap:react-modele} (le tour de React, React-tRace et sa lecture réduite de la boucle) ; \cref{chap:interpretation-abstraite} (\cref{def:point-fixe}, \cref{def:widening}, \cref{def:soundness}) ; \cref{chap:statevalue-stabilite} (\cref{def:stabilite}) ; \cref{chap:transfert-stores} (environnements, \cref{def:store}, \code{exec_setter_call}, \code{recompute_memo}) ; \cref{chap:cfg-analyzer-dominance} (\code{analyze_cfg}, sa passe et la monotonie du store qu'elle rend) ; \cref{chap:splice-inlining} pour la mécanique des greffes.}

Le \cref{chap:cfg-analyzer-dominance} sait exécuter abstraitement \emph{un}
corps : le rendu, un effet, un handler. Mais un composant React n'est pas un
corps, c'est un cycle : le rendu produit un environnement, les effets le lisent
et écrivent l'état, l'état modifié provoque un nouveau rendu, et ainsi de suite
jusqu'à ce que plus rien ne change --- ou indéfiniment, ce qui est précisément
l'un des défauts à détecter. Ce chapitre décrit la fonction qui referme ce
cycle, \code{analyze_component_impl} (\src{src/engine/fixpoint.rs}{130}{698}),
et tout ce qu'elle appelle. C'est la pièce centrale du moteur : elle reçoit un
\code{ComponentIR} et rend un \code{AnalysisResult}, l'objet que toutes les
règles lisent.

Le \cref{chap:react-modele} en a donné, sous forme d'algorithme, une
lecture réduite, suffisante pour la correspondance avec React-tRace. Ici nous
la reprenons à la taille du code, ligne à ligne : ce qui se passe une fois
avant la boucle, ce que fait chacune des quatre passes d'un tour, comment on
décide qu'on a fini, ce qu'on recalcule ensuite et pourquoi. La progression va
du composant le plus simple (un \code{useState}, un effet) aux cas qui ont
imposé les choix les moins évidents : handlers dans la boucle, mémos hors
du test de convergence, ré-exécution des effets depuis $\Bot$. Le chapitre
se termine par un inventaire honnête de ce que le moteur garantit et de ce
qu'il ne garantit pas, au commit \snapshotCommit.

Toutes les sorties d'analyseur ont été observées avec le binaire construit à
ce commit sur des fichiers placés sous \file{/tmp/ch13/}, par
\code{reactant check --verbose --info --no-color --project plain <fichier>}
(la ligne \code{[verbose] X: N iteration(s), widened: [...]} vient de
\src{src/driver/mod.rs}{419}{428} et affiche \code{iterations} et les clés de
\code{widen_trace}). Les valeurs internes des stores, invisibles depuis la
ligne de commande, ont été lues avec une petite sonde Rust hors dépôt qui lie
la crate \code{reactant} par chemin et imprime les champs de
l'\code{AnalysisResult} ; sa notation est \code{s0=num[0,3]} pour \og le slot
0 vaut l'intervalle $[0,3]$\fg{}, \code{ref:PerRender} pour une référence
fraîche à chaque rendu, \code{⊤} pour \code{StateValue::top()}.

% ===========================================================================
\section{Le problème : un ensemble d'états, pas une simulation}
\label{sec:point-fixe-composant-probleme}

\subsection{Un compteur qui ne s'arrête jamais}

Le composant le plus court qui boucle à l'infini est celui-ci :

\begin{lstlisting}[language=TypeScript,caption={La boucle infinie canonique (\file{/tmp/ch13/runaway.tsx})},label={lst:point-fixe-composant-runaway}]
import { useState, useEffect } from "react";

export function Runaway() {
  const [x, setX] = useState(0);
  useEffect(() => setX(x + 1));
  return <div>{x}</div>;
}
\end{lstlisting}

Concrètement (un tour de React au sens du \cref{chap:react-modele}) : React rend \code{Runaway} avec
$x = 0$, commit, exécute l'effet, qui appelle \code{setX(1)} ; l'état ayant
changé, React rend à nouveau avec $x = 1$, l'effet tourne encore (il n'a pas
de tableau de deps), \code{setX(2)}, et ainsi de suite. La suite des valeurs
de \code{x} est $0, 1, 2, 3, \dots$ ; le composant ne se stabilise jamais.
L'analyseur le dit :

\begin{sortie}
$ reactant check --verbose --info --no-color --project plain runaway.tsx
[verbose] symbol graph: 1 nodes, topo order = [Runaway@runaway.tsx]
[verbose] 1 components analyzed
[verbose] cache hits: 0, misses: 0
  [verbose] Runaway: 3 iteration(s), widened: [0]
  Runaway  (2 hooks)  runaway.tsx
    warn   infinite-loop  [hook:0]  (line 5:2)  this effect keeps pushing state `x` to new values on every run. Potential infinite render loop
       (2 trace step(s), rerun with --trace)
    info   widening-info  (line 5:2)  state `x` kept changing during analysis and was approximated to converge, so findings that depend on it may be imprecise
       (2 trace step(s), rerun with --trace)
    verified  conditional-hook  all hooks run unconditionally, in a stable order
    [... six autres lignes verified ...]

⚠  1 warning(s) across 1 file(s).
\end{sortie}

Comment un programme qui ne fait que \emph{trois} itérations peut-il conclure
quelque chose sur une exécution qui n'en finit pas ? La réponse est le sujet du
chapitre, mais l'idée tient en deux phrases. Le moteur ne suit pas la suite
$0, 1, 2, \dots$ ; il calcule, pour le slot de \code{x}, un \emph{ensemble}
de valeurs qui contient toutes celles que \code{x} peut prendre, à n'importe
quel tour. Cet ensemble est représenté par une valeur abstraite (ici un
intervalle), il grossit tant que l'effet écrit quelque chose de neuf, et
quand il grossit trop longtemps, un opérateur de widening
(\cref{def:widening}) le fait sauter à $[0, +\infty]$ : \og \code{x} peut
prendre n'importe quelle valeur $\geq 0$\fg{}. Ce saut est le signal de
divergence que la règle \regle{infinite-loop} lit.

\subsection{Pourquoi pas simuler}

On pourrait vouloir exécuter le composant : rendre, appliquer les effets,
rendre à nouveau, et déclarer une boucle infinie au bout de, disons, cent
tours. Trois raisons rendent cette voie impraticable pour un analyseur
statique.

\begin{enumerate}
  \item \emph{Les entrées sont inconnues.} Un composant reçoit des props, lit
    des contextes, appelle des bibliothèques. Aucune valeur concrète n'est
    disponible ; il faut raisonner sur \emph{toutes} les valeurs possibles à la
    fois, c'est-à-dire sur des ensembles.
  \item \emph{Les événements sont arbitraires.} Un \code{onClick} s'exécute
    zéro, une ou mille fois, dans n'importe quel ordre par rapport aux autres
    handlers. Aucune suite finie de tours ne représente toutes les
    exécutions.
  \item \emph{La terminaison n'est pas décidable}, et surtout, on veut une
    réponse \emph{sound} (\cref{def:soundness}) : si l'analyseur dit
    \og rien à signaler\fg, aucune exécution concrète ne doit boucler. Une
    simulation bornée peut manquer une boucle qui n'apparaît qu'à la
    cent-unième itération ou qu'avec une entrée particulière.
\end{enumerate}

L'interprétation abstraite répond aux trois points d'un coup : on remplace
chaque valeur par un élément d'un treillis qui représente un ensemble de
valeurs, on remplace la suite des tours par une suite \emph{croissante}
d'états abstraits, et on s'arrête à un post-point fixe (\cref{def:point-fixe})
qui, par le théorème du post-point fixe du
\cref{chap:interpretation-abstraite}, contient tout état concret atteignable.

\begin{definition}{État abstrait d'un composant, tour abstrait}{point-fixe-composant-etat}
L'\emph{état abstrait} d'un composant est son \code{StateStore}
(\cref{def:store}) : une fonction $S$ des labels de ses hooks \code{useState}
vers des \code{StateValue}, un label absent valant $\Bot$. Un \emph{tour
abstrait} est une exécution de la boucle de \code{analyze_component_impl} :
une passe de rendu depuis $S$, le recalcul des mémos, une passe par effet
et une passe par handler, toutes depuis $S$, puis le join $N$ des stores
produits. La \emph{suite des tours} est $S_0 \lleq S_1 \lleq \dots$, où $S_0$
est le store amorcé par les initialiseurs et $S_{k+1}$ vaut $N_k$ ou
$S_k \widen N_k$. La boucle s'arrête au premier $k$ tel que
$N_k \lleq S_k$ ; ce $S_k$ est l'état abstrait convergé.
\end{definition}

La \cref{fig:point-fixe-composant-chronogramme} met en regard le temps de
React et celui de l'analyseur pour \code{Runaway}. Le point important est que
le tour abstrait $k$ ne représente pas le $k$-ième tour de React : il
représente \emph{tous} les états atteignables en au plus $k$ tours, et un
tour de plus, quand il n'ajoute rien, prouve que l'ensemble est clos.

\begin{figure}[htbp]\centering
\resizebox{\linewidth}{!}{%
\begin{tikzpicture}[font=\scriptsize,node distance=2mm]
  % --- ligne React (concret)
  \node[anchor=east] at (-0.9,0) {\textbf{React}};
  \foreach \i/\v/\w in {0/0/1,1/1/2,2/2/3,3/3/4} {
    \node[bloc,minimum width=11mm,inner sep=2pt] (r\i) at ({\i*3.8},0) {rendu\\$x=\v$};
    \node[bloc,minimum width=9mm,inner sep=2pt,fill=white,right=of r\i] (e\i) {effet\\\texttt{setX(\w)}};
  }
  \foreach \i/\j in {0/1,1/2,2/3} { \draw[fleche] (e\i) -- node[etiquette,above]{commit} (r\j); }
  \foreach \i in {0,1,2,3} { \draw[fleche] (r\i) -- (e\i); }
  \node[right=of e3] {$\cdots$};
  % --- ligne abstraite
  \node[anchor=east] at (-0.9,-2.3) {\textbf{moteur}};
  \node[bloc,minimum width=13mm,inner sep=2pt,draw=ra-amber] (s0) at (0,-2.3) {$S_0$\\$[0,0]$};
  \node[bloc,minimum width=13mm,inner sep=2pt,draw=ra-amber] (s1) at (3.8,-2.3) {$S_1$\\$[0,1]$};
  \node[bloc,minimum width=13mm,inner sep=2pt,draw=ra-amber] (s2) at (7.6,-2.3) {$S_2$\\$[0,2]$};
  \node[bloc,minimum width=13mm,inner sep=2pt,draw=ra-amber] (s3) at (11.4,-2.3) {$S_3$\\$[0,+\infty]$};
  \node[bloc,minimum width=13mm,inner sep=2pt,draw=ra-amber,fill=white] (s4) at (15.2,-2.3) {$N_3 \lleq S_3$\\fin};
  \draw[fleche] (s0) -- node[etiquette,above]{tour 0 : join} (s1);
  \draw[fleche] (s1) -- node[etiquette,above]{tour 1 : join} (s2);
  \draw[fleche] (s2) -- node[etiquette,above,align=center]{tour 2 : $\widen$\\(\code{iteration} = 3)} (s3);
  \draw[fleche] (s3) -- node[etiquette,above]{tour 3} (s4);
  % --- liens de concrétisation
  \foreach \i in {0,1,2} { \draw[flechemay] (r\i.south) -- node[etiquette,right,text=ra-amber]{$\in \gam$} (s\i.north); }
  \draw[flechemay] (r3.south) -- node[etiquette,right,text=ra-amber]{$\in \gam$} (s3.north);
  \node[etiquette,align=left,anchor=west,text width=150mm] at (-2.2,-3.7) {Chaque rendu concret de React lit une valeur de \code{x} contenue dans la concrétisation $\gam(S_k)$ du store abstrait de même rang ; le store $S_3$ contient aussi tous les rendus suivants.};
\end{tikzpicture}}
\caption{Chronogramme de \code{Runaway} : les tours de React (en haut,
infinis) et les tours de la boucle externe (en bas, quatre). Le store abstrait
$S_k$ sur-approxime toutes les valeurs de \code{x} atteignables en au plus $k$
tours ; après le widening, $S_3$ les contient toutes.}
\label{fig:point-fixe-composant-chronogramme}
\end{figure}

\subsection{Ce qu'un ensemble ne dit pas}

Le prix de cette abstraction est visible sur un second exemple. Ici l'état
est un objet, recréé par l'effet qui en dépend :

\begin{lstlisting}[language=TypeScript,caption={Un objet recréé dans ses propres deps (\file{/tmp/ch13/fresh.tsx})},label={lst:point-fixe-composant-fresh}]
import { useState, useEffect } from "react";

export function Fresh() {
  const [o, setO] = useState({ n: 0 });
  useEffect(() => {
    setO({ n: o.n + 1 });
  }, [o]);
  return <div>{o.n}</div>;
}
\end{lstlisting}

Concrètement, c'est une boucle infinie : chaque exécution de l'effet stocke
un objet neuf, \code{Object.is} échoue sur le dep \code{[o]}, l'effet
tourne encore. Abstraitement, la valeur du slot est \code{ref:PerRender}
(\cref{def:stabilite} : \og une référence fraîche à chaque rendu\fg) dès
l'amorçage, l'effet écrit \code{ref:PerRender}, et
$\mathit{PerRender} \join \mathit{PerRender} = \mathit{PerRender}$ : le
store ne bouge pas, la boucle s'arrête \emph{au premier tour}, sans
widening :

\begin{sortie}
  [verbose] Fresh: 0 iteration(s), widened: []
  Fresh  (2 hooks)  fresh.tsx
    error  infinite-loop  [hook:0]  (line 5:2)  this effect recreates object state `o` it depends on. Every run stores a fresh reference (`Object.is` always fails) and re-triggers itself: infinite render loop
\end{sortie}

L'\Error{} est pourtant là, et c'est une \Error{}, pas un \Warning{} : elle
ne vient pas du point fixe mais des \emph{relations} dérivées après lui
(\cref{chap:relations}, \cref{chap:churn}) --- ici la ligne \code{exact =
true} de \relation{effect_triggers} et la fraîcheur de l'écriture dans
\relation{slot_writers}. La leçon, que le \cref{chap:relations} reprend à son
compte (\adr{42}), est celle-ci : le point fixe calcule un \emph{ensemble de
valeurs}, jamais une \emph{suite de transitions}. La suite $0 \to 1 \to 0 \to
1$ et le fait \og atteint $\{0, 1\}$\fg{} ont exactement la même abstraction.
Une question qui porte sur les transitions (\og cette écriture
relance-t-elle cet effet ?\fg) n'a pas de réponse dans le store ; elle en a
une dans les relations, calculées à convergence sur les CFG. Ce chapitre
s'arrête donc à la frontière : il montre comment le store converge et quelles
données sont figées pour les relations, et laisse leur calcul aux chapitres
suivants.

\subsection{Les quatre portes d'entrée et la configuration}

Le fichier \file{src/engine/fixpoint.rs} (2\,815 lignes, dont un millier de
tests) expose quatre fonctions publiques, réexportées par
\file{src/engine/mod.rs}, qui aboutissent toutes au cœur privé
\code{analyze_component_impl}.

\rustsnippet[label={lst:point-fixe-composant-entrees}]{src/engine/fixpoint.rs}{83}{118}{les deux entrées intra-composant}

\begin{itemize}
  \item \code{analyze_program} (\srcline{src/engine/fixpoint.rs}{704}) est
    l'entrée de production, appelée par le driver : elle analyse les racines
    avec un \code{InterCtx} (phase 1), puis les composants non atteints sans
    (phase 2). Le \cref{sec:point-fixe-composant-programme} en dit ce qui
    concerne la boucle ; le \cref{chap:inter-composants} fait le reste.
  \item \code{analyze_component} et \code{analyze_component_as}
    (\cref{lst:point-fixe-composant-entrees}) analysent un composant
    \emph{isolé} : environnement d'entrée $\Bot$ (toute variable inconnue vaut
    $\Top$, \cref{chap:transfert-stores}), tas neuf, \code{inter = None}. La première
    l'identifie par \code{ComponentId::SYNTHETIC}
    (\srcline{src/ir/component_id.rs}{38}), la seconde par un identifiant déjà
    interné. Ce sont les entrées de la plupart des tests unitaires et de
    plusieurs tests d'intégration (\file{tests/widening_e2e.rs},
    \file{tests/functional_updater.rs}).
  \item \code{analyze_component_inter} (\src{src/engine/fixpoint.rs}{65}{81})
    est la fonction de rappel que \code{eval_comp_app} appelle pour analyser
    un enfant rencontré dans un rendu ; elle existe pour casser la dépendance
    circulaire entre \code{domains::transfer} et \code{engine::fixpoint}
    (\cref{chap:transfert-stores}).
\end{itemize}

\begin{piege}
Une analyse \emph{intra} (\code{analyze_component}) et une analyse
\emph{inter} (\code{analyze_program}) du même composant ne calculent pas la
même chose. En intra, il n'y a pas de \code{HookRegistry} donc pas d'inlining
des hooks personnalisés (\cref{sec:point-fixe-composant-expansion}), pas de
\code{SharedStateStore} (\cref{sec:point-fixe-composant-join}), et un
\code{SummaryRegistry} vide par défaut. Un test qui passe par
\code{analyze_component} ne prouve rien du chemin suivi par la ligne de
commande ; l'\issue{14} relève que 36 des 42 fichiers de tests d'intégration
construisent \code{Config::default()}, \og which no user ever runs\fg. Nous
en verrons une conséquence concrète au
\cref{sec:point-fixe-composant-trous} (l'updater fonctionnel).
\end{piege}

Les paramètres du moteur tiennent dans une structure :

\rustsnippet[label={lst:point-fixe-composant-config}]{src/engine/fixpoint.rs}{38}{60}{les paramètres du moteur et leurs défauts}

\code{widen_threshold} sert \emph{deux fois} : c'est le nombre de tours
externes qui changent l'état avant que la boucle passe du join au widening
(\cref{sec:point-fixe-composant-convergence}), et c'est aussi le nombre de
passages par un arc arrière avant le widening d'un en-tête de boucle dans
\code{analyze_cfg} (\cref{chap:cfg-analyzer-dominance}). Les deux registres
alimentent l'expansion (\cref{sec:point-fixe-composant-expansion}) ;
\code{max_inline_depth} borne le nombre de greffes d'utilitaires par CFG. Le
driver n'utilise pas \code{Config::default()} tel quel : il remplace le
\code{SummaryRegistry} vide par \code{SummaryRegistry::new_with_common()}
(\src{src/driver/mod.rs}{370}{374}), qui connaît les hooks des bibliothèques
courantes.

La signature du cœur, avec le plan que son auteur en a écrit :

\rustsnippet[label={lst:point-fixe-composant-signature}]{src/engine/fixpoint.rs}{120}{138}{la doc d'en-tête et la signature de \code{analyze_component_impl}}

Un détail à corriger d'emblée : l'étape 1 du plan (\og Import cross-component
state\fg) n'a pas lieu en tête de boucle. L'import de la tranche du
\code{SharedStateStore} se fait dans le test de convergence
(\src{src/engine/fixpoint.rs}{488}{493}), après les passes ; nous y
reviendrons. Le reste du plan est fidèle, et la
\cref{fig:point-fixe-composant-flot} le dessine en entier, préparation et
post-convergence comprises. Les trois sections suivantes en parcourent les
trois zones.

\begin{figure}[htbp]\centering
\begin{tikzpicture}[font=\scriptsize,node distance=3mm and 6mm,
    etape/.style={bloc,text width=50mm,align=left,inner sep=3pt},
    zone/.style={draw=ra-gray,dashed,rounded corners=3pt,inner sep=4pt}]
  % --- préparation
  \node[etape] (consts) {constantes de module $\to$ \code{initial_env}, \code{module_env} \hfill l.~153--190};
  \node[etape,below=of consts] (util) {\code{expand_utility_calls} (greffes d'utilitaires) \hfill l.~202--218};
  \node[etape,below=of util] (hooks) {\code{expand_custom_hooks} (greffes de hooks, résumés) \hfill l.~221--229};
  \node[etape,below=of hooks] (thr) {\code{custom_arg_returns} ; seuils $T$ ; \code{callback_bodies} \hfill l.~240--295};
  \node[etape,below=of thr] (seed) {amorçage : $S[\ell] \join= \sem{\mathit{init}}$ pour chaque \code{State} \hfill l.~314--353};
  \begin{scope}[on background layer]
    \node[zone,fit=(consts)(seed),label={[etiquette,anchor=south west]north west:préparation (une fois)}] (prep) {};
  \end{scope}
  % --- boucle
  \node[etape,right=of consts,xshift=4mm] (render) {rendu : \code{analyze_cfg(render_cfg, initial_env, S)} $\to$ \code{block_states}, $R$ \hfill l.~355--380};
  \node[etape,below=of render] (memo) {\code{env_exit} ; mémos : \code{recompute_memo(deps, env_exit)} \hfill l.~382--413};
  \node[etape,below=of memo] (eff) {effets : $E = \bigsqcup$ \code{analyze_cfg(body, env_exit, S)} \hfill l.~415--449};
  \node[etape,below=of eff] (hdl) {handlers : $H = \bigsqcup$ \code{analyze_cfg(body, env_exit, S)} \hfill l.~451--483};
  \node[etape,below=of hdl] (join) {$N = R \join E \join H \join$ \code{shared_state.slice(id)} \hfill l.~485--493};
  \node[etape,below=of join,draw=ra-amber] (test) {$N \lleq S$ ? \hfill l.~495};
  \node[etape,below=of test] (widen) {\code{iteration}++ ; $S := N$ ou $S := S \widen_T N$ (+ \code{widen_trace}) ; cap 100 \hfill l.~499--529};
  \begin{scope}[on background layer]
    \node[zone,fit=(render)(widen),label={[etiquette,anchor=south west]north west:boucle externe (\code{loop}, l.~355--530)}] (loop) {};
  \end{scope}
  % --- post
  \node[etape,below=of widen,yshift=-6mm] (refresh) {rafraîchissement : \code{analyze_cfg(render_cfg, initial_env, S)}, \code{inter = None} \hfill l.~532--563};
  \node[etape,below=of refresh] (pure) {\code{effect_setter_writes} : effets rejoués depuis $\Bot$ \hfill l.~565--594};
  \node[etape,below=of pure] (rel) {relations : \relation{slot_writers}, \relation{slot_seeds}, \relation{registrations}, \relation{effect_triggers} \hfill l.~596--667};
  \node[etape,below=of rel,fill=white] (out) {\texttt{AnalysisResult \{ state\_store: S, iterations, widen\_trace, \dots \}} \hfill l.~670--697};
  \begin{scope}[on background layer]
    \node[zone,fit=(refresh)(out),label={[etiquette,anchor=south west]north west:post-convergence}] (post) {};
  \end{scope}
  % --- flèches
  \foreach \a/\b in {consts/util,util/hooks,hooks/thr,thr/seed,render/memo,memo/eff,eff/hdl,hdl/join,join/test,refresh/pure,pure/rel,rel/out} { \draw[fleche] (\a) -- (\b); }
  \draw[fleche] (seed.east) -- ++(4mm,0) |- (render.west);
  \draw[fleche] (test) -- node[etiquette,right]{non} (widen);
  \draw[fleche] (widen.east) -- ++(4mm,0) |- node[etiquette,pos=0.25,right]{tour suivant} (render.east);
  \draw[flechemust] (test.west) -- ++(-4mm,0) |- node[etiquette,pos=0.25,left]{oui} (refresh.west);
\end{tikzpicture}
\caption{Le flot de \code{analyze_component_impl}
(\file{src/engine/fixpoint.rs}). Les numéros de ligne sont ceux du commit
\snapshotCommit. La préparation s'exécute une fois ; la boucle externe
répète les quatre passes jusqu'à $N \lleq S$ ; la post-convergence recalcule
ce que les règles liront.}
\label{fig:point-fixe-composant-flot}
\end{figure}

% ===========================================================================
\section{La préparation : ce qui se passe une fois}
\label{sec:point-fixe-composant-preparation}

Avant le premier tour, \code{analyze_component_impl} transforme le
\code{ComponentIR} qu'il reçoit (\src{src/engine/fixpoint.rs}{139}{353}).
Rien de ce qui suit n'est répété : les CFG sont modifiés en place par les
greffes, les seuils sont récoltés une fois, l'état est amorcé une fois. La
boucle n'a ensuite plus qu'à exécuter.

\subsection{Les constantes de module}

Un composant lit souvent des constantes définies au niveau du module
(\code{const LIMIT = 10;}, \code{const EMPTY = [];},
\code{const Ctx = createContext(...)}). Sans traitement particulier, une
lecture de \code{LIMIT} dans le rendu tomberait sur une variable absente de
l'environnement, donc sur $\Top$ (\cref{chap:transfert-stores}) --- et une constante
\code{EMPTY} passée en dep se lirait \og peut changer à chaque rendu\fg.

\rustsnippet[label={lst:point-fixe-composant-consts}]{src/engine/fixpoint.rs}{153}{190}{les constantes de module amorcent deux environnements}

Le lowering a rangé ces constantes dans \code{comp.module_consts}
(\cref{chap:ir}), sous la forme d'un \code{ModuleConstInit} : un littéral
primitif (\code{Prim}), une référence (\code{Ref}) ou un objet de contexte
(\code{Context}). La boucle les évalue une fois, contre des stores vides et un
environnement vide, et les place dans deux environnements distincts :
\begin{itemize}
  \item \code{initial_env}, l'environnement d'entrée du rendu, où une liaison
    déjà présente --- une prop liée par le parent en analyse inter --- l'emporte
    sur la constante de même nom ;
  \item \code{module_env}, qui ne contient \emph{que} les constantes, et qui
    servira d'environnement d'entrée aux corps des arguments de hooks opaques
    (voir plus bas) : une constante de module a la même valeur en tout point
    du programme, c'est la seule liaison que l'on puisse y semer sans risque.
\end{itemize}

\begin{react}[Constantes de module]
Un module ECMAScript est évalué une fois par vie du programme ; une
\code{const} de module ne change donc jamais entre deux rendus, quel que soit
le composant qui la lit. C'est ce qui autorise \code{Stability::Stable} pour
une référence de module, alors qu'un littéral objet écrit \emph{dans} le corps
du composant vaut \code{PerRender} (\cref{def:stabilite}). Un objet
\code{createContext(...)} est une référence de module comme une autre ; le
rôle \code{Context} qu'il porte en plus n'est lu que par les règles
(\code{unstable-context-value}), pas par le moteur.
\end{react}

\subsection{L'expansion : les utilitaires, puis les hooks personnalisés}
\label{sec:point-fixe-composant-expansion}

Le composant que reçoit la boucle n'est pas toujours celui que le programmeur
a écrit. Un \code{useTicker(5)} n'est pas un hook natif : l'IR le représente
par une entrée \code{HookEntry::Custom} et un \code{HookMarker} au site
d'appel, et rien dans le domaine ne sait ce qu'il fait. De même, un appel
\code{helper(x);} en position d'instruction est un \code{Call} opaque, qui
s'évalue à $\Top$. La préparation remplace ce qu'elle peut par le corps
correspondant, greffé dans le CFG de l'appelant --- c'est l'\emph{inlining}
du \cref{chap:splice-inlining}, dont nous ne reprenons ici que ce que le point
fixe en fait.

\rustsnippet[label={lst:point-fixe-composant-expansion}]{src/engine/fixpoint.rs}{192}{229}{les deux expansions, dans cet ordre}

L'ordre est voulu (commentaire des l.~200--201) : les utilitaires d'abord,
pour qu'un hook appelé \emph{dans} un utilitaire devienne visible à
l'expansion des hooks ; les hooks ensuite, \emph{avant} l'amorçage, pour que
les \code{useState} d'un hook greffé soient amorcés comme les autres.

\paragraph{Un exemple.} Le composant \code{Clock} appelle un hook
personnalisé qui, à lui seul, boucle :

\begin{lstlisting}[language=TypeScript,caption={Un hook personnalisé qui boucle (\file{/tmp/ch13/ticker.tsx})},label={lst:point-fixe-composant-ticker}]
import { useState, useEffect } from "react";

function useTicker(start: number) {
  const [t, setT] = useState(start);
  useEffect(() => {
    setT(t + 1);
  });
  return t;
}

export function Clock() {
  const t = useTicker(5);
  return <div>{t}</div>;
}
\end{lstlisting}

La sonde montre les deux visages du même composant. En \emph{intra}
(\code{analyze_component}), l'entrée \code{hook 0: Custom useTicker} n'est
pas expansée ; le store est vide et la boucle s'arrête au premier tour
(\code{iterations = 0}). En \emph{inter} (\code{analyze_program}, le chemin
du CLI), l'entrée \code{Custom} de label 0 a été remplacée par un
\code{State} de label 1 et un \code{Effect} de label 2, l'init \code{start} a
été substitué par \code{5}, et le point fixe trouve la divergence :

\begin{sortie}
## engine result Clock
  iterations = 3
  state_store = s1=num[5,inf]
  effect_setter_writes = s1=num[6,inf]
  widen_trace = [(1, WidenEvent { iteration: 3, writers: [2] })]
\end{sortie}
\begin{sortie}
  [verbose] Clock: 3 iteration(s), widened: [1]
  Clock  (1 hooks)  ticker.tsx
    warn   infinite-loop  [hook:1]  (line 5:2)  this effect keeps pushing state `t` to new values on every run. Potential infinite render loop
\end{sortie}

Le rapport dit \code{(1 hooks)} --- le composant tel qu'écrit n'a qu'un
appel de hook --- mais désigne \code{[hook:1]}, le label du \code{useState}
greffé, et pointe la ligne 5, dans le corps de \code{useTicker}. C'est la
provenance (\code{inline_origins}, \adr{19}) qui permet aux règles de
retrouver ce fichier et cette ligne.

\paragraph{Ce que fait \code{expand_custom_hooks}.} La fonction
(\src{src/engine/fixpoint.rs}{889}{1142}) parcourt \code{hooks} par indice et,
pour chaque entrée \code{Custom} qu'elle sait résoudre, la remplace par les
sous-entrées du hook. Les étapes, avec leurs lignes :

\begin{enumerate}
  \item \emph{Sortie immédiate} sans \code{inter} ou sans
    \code{HookRegistry} (l.~898--901) : c'est pourquoi l'analyse intra ne
    greffe rien.
  \item \emph{Garde de récursion par nom} (l.~927--931) : un ensemble
    \code{expanding} des noms déjà greffés ; une seconde occurrence du même
    nom est sautée et reste opaque.
  \item \emph{Résolution} (l.~933--942) : d'abord la clé
    \code{(fichier résolu, nom)}, sinon le premier hook de ce nom, trié par
    clé (\code{get_by_name}). C'est plus permissif que la résolution des
    utilitaires, qui est fail-closed (voir plus bas).
  \item \emph{Repli sur un résumé} (l.~943--994) : un hook absent du registre
    mais connu du \code{SummaryRegistry} (TanStack, React Router\dots) n'est
    pas greffé ; son \code{HookMarker} est ré-étiqueté avec un
    \code{SummaryValue} (\code{retag_marker}), et l'entrée \code{Custom} est
    \emph{conservée} pour que les règles des hooks voient toujours le site
    d'appel.
  \item \emph{Décalage des labels} (l.~996--997) : \code{offset} est le
    premier label libre ; les sous-entrées du hook sont renumérotées
    à partir de lui, ce qui donne le label 1 au \code{useState} de
    \code{Clock}.
  \item \emph{Substitution des paramètres dans les inits}
    (l.~999--1007) : les initialiseurs de \code{useState} sont évalués dans
    un environnement séparé qui ne voit pas les liaisons du rendu
    (\cref{sec:point-fixe-composant-amorcage}) ; \code{useState(start)} doit
    donc devenir \code{useState(5)} par substitution syntaxique.
  \item \emph{Sel et identités} (l.~1012--1026) : \code{salt.take(span)}
    fournit un sel d'alpha-renommage (les locales du hook deviennent
    \code{t\#0}, \code{setT\#0}) et un décalage d'\code{ExprId} pour les
    sites d'allocation ; la table de renommage évite les liaisons de
    l'appelant (\issue{141}).
  \item \emph{Greffe du corps entier} au site du \code{HookMarker}
    (l.~1043--1078), retour lié à la variable de l'appelant, et
    enregistrement de la région greffée (\adr{27} §4 : une écriture dans
    cette région est \og via \code{useTicker}\fg, pas directe). À défaut de
    marqueur, cas défensif, seules les instructions du bloc d'entrée sont
    greffées et la provenance du rendu est empoisonnée
    (\code{render_poisoned}, fail-closed, l.~1079--1109).
  \item \emph{Provenance} (l.~1111--1130) : un \code{InlineOrigin} et les
    lignes \code{HookProvenance} du hook, décalées et marquées
    \code{inlined: true} (\adr{23}).
  \item \emph{Remplacement sans avancer} (l.~1132--1141) : l'entrée
    \code{Custom} est remplacée par les sous-entrées, et l'indice \code{i}
    n'est pas incrémenté, car la première sous-entrée peut elle-même être un
    \code{Custom} : l'expansion est récursive, bornée par la garde de
    récursion.
\end{enumerate}

Les deux réserves d'identités que consomme chaque greffe vivent dans une
petite structure :

\rustsnippet[label={lst:point-fixe-composant-spliceids}]{src/engine/fixpoint.rs}{828}{867}{les identités d'une greffe : sel et curseur d'\code{ExprId}}

\begin{invariant}{Réserves monotones sur tout le composant}{point-fixe-composant-spliceids}
Le sel et le curseur d'allocation de \code{SpliceIds} ne font que croître
pendant toute la préparation d'un composant, greffes d'utilitaires, greffes
de hooks et formes de résumé (\code{alloc_one}) comprises. Deux greffes
n'occupent donc jamais le même site d'allocation, et le même hook greffé
deux fois ne partage pas une entrée du tas (\issue{134}).
\end{invariant}

\begin{decision}[ADR-005 --- portée intra-procédurale et registre modulaire de hooks]
\emph{Décision} (\adr{5}, 2026-05-29). Chaque composant est analysé
indépendamment ; un hook inconnu rend \code{Unknown} ; un registre de
modèles de hooks à quatre couches (natifs, bibliothèques, configuration
utilisateur, repli) évite l'inlining ; une phase 2 ultérieure ajoutera un
inlining \emph{call-string-1}. \emph{Ce qu'il en est au commit
\snapshotCommit.} L'inlining des hooks personnalisés existe
(\code{expand_custom_hooks}), et il n'est pas limité à la profondeur 1 : il
est récursif, gardé par nom. Le trait \code{HookModel} de l'ADR est devenu
\code{HookSummary} (\file{src/registry/summary.rs}) et la configuration est
\file{reactant.config.json}, pas \file{reactant.toml}. L'ADR n'a pas été
amendé ; sa lecture doit se faire avec ces trois corrections.
\end{decision}

\paragraph{Les utilitaires.} \code{expand_utility_calls}
(\src{src/engine/fixpoint.rs}{1554}{1602}) parcourt le rendu puis chaque corps
d'\code{Effect}, \code{Memo}, \code{Callback} et \code{Handler} --- pas les
inits de \code{State} ni les arguments des \code{Custom} --- et appelle
\code{inline_in_cfg} sur chacun :

\rustsnippet[label={lst:point-fixe-composant-inline}]{src/engine/fixpoint.rs}{1611}{1634}{la boucle de greffes d'un CFG et son budget}

Trois propriétés à retenir. La cible est une \emph{instruction}
(\code{let x = util(...)} ou \code{util(...);}, \code{utility_call_target},
l.~1698--1718) : un appel en position d'expression reste opaque
(\issue{52}). La résolution du nom est fail-closed
(\code{FunctionRegistry::resolve}, \src{src/engine/function_registry.rs}{51}{63}) :
une définition dans le fichier de la région où se trouve l'appel, sinon
l'arête d'import résolue au lowering pour ce nom local, jamais une devinette
par nom à travers les fichiers. Enfin le budget \code{max_inline_depth} est
\emph{par CFG} (\code{let mut budget = ctx.max_depth} au début de chaque
\code{inline_in_cfg}), décrémenté même quand la greffe échoue, et chaque nom
n'est greffé qu'une fois par CFG (\code{expanding}, neuf pour chaque CFG) :
un second appel du même utilitaire dans le même corps reste un \code{Call}
opaque.

\begin{piege}
Le drapeau \code{truncated}, qui signale qu'un appel est resté opaque faute
de budget, n'est enregistré que sous \code{Some(inter)}
(\cref{lst:point-fixe-composant-expansion}, l.~211) : il vit dans
\code{AnalysisStats}, accessible seulement par l'\code{InterCtx}. En
phase 2 --- et en intra --- il est jeté. Nous l'avons vérifié avec un
composant \code{A} qui appelle dix utilitaires \code{u1(1); ... u10(1);}
(budget 8), dans deux versions. Quand \code{A} est une racine (il rend un
\code{<div>}), le CLI émet
\code{info analysis-limit utility inlining ran out of splice budget here, so the remaining utility calls are treated as unknown (FN possible) [...]}
puis \code{suspended analysis-limit 9 passing check(s) withheld}. Quand
\code{A} et \code{B} se rendent mutuellement --- aucune racine, tous deux en
phase 2 ---, la troncature a lieu de même, mais le CLI publie pour \code{A}
neuf lignes \code{verified} et aucune Info. Les assurances couvrent alors
du code que l'analyse n'a pas lu, contrairement à l'intention écrite au
commentaire de \code{inline_in_cfg} (l.~1621--1625). Ce n'est pas une
sous-approximation du store (les appels opaques valent $\Top$), mais un
défaut de la couche de certification ; aucune issue ne le décrit au commit
\snapshotCommit. Le correctif naturel, au niveau central, serait de porter le
drapeau sur l'\code{AnalysisResult} plutôt que sur les statistiques inter.
\end{piege}

\begin{piege}
La garde de récursion par \emph{nom} coupe aussi le second appel légitime
d'un même hook. Avec \code{const a = useTicker(5); const b = useTicker(9);},
seul le premier est greffé ; pour le second le CLI émet
\code{info analysis-limit [hook:1] (line 13:8) hook `useTicker` was not found in the registry. Pass its source file or add a HookSummary to analyse it (FN possible)}
puis suspend sept assurances. Le comportement est sound (le second appel
vaut $\Top$ et le composant retient ses assurances), mais le message est
trompeur : le hook \emph{est} dans le registre. Le commentaire de la garde
(l.~884--888) assume ce cas comme \og rare\fg{} et \og FN, not FP\fg{} ;
c'est un compromis, pas un défaut, au sens de \file{docs/limitations.md}.
\end{piege}

\paragraph{Les arguments des hooks restés opaques.} Un \code{useStore(s => s.count)}
qui n'a été ni greffé ni résumé garde son entrée \code{Custom}. Le moteur
exécute tout de même le corps de ses arguments \code{FnLit} (ou d'une variable
dont la liaison à un littéral de fonction est certifiée syntaxiquement,
\code{certified_fn_binding}), avec ses paramètres liés à $\Top$ sur
\code{module_env}, et range la valeur de retour jointe dans
\code{custom_arg_returns}, clé \code{(label, indice)}
(\src{src/engine/fixpoint.rs}{231}{277}). L'argument de soundness tient en
une phrase du commentaire : une capture non constante lit $\Top$, qui
sur-approxime sa valeur en tout point du programme, donc quel que soit le
moment où la bibliothèque appelle le sélecteur, la valeur stockée le couvre.
Ce champ sert au vocabulaire déclaratif (\adr{23} §3, \cref{chap:regles-declaratives}).

\subsection{Les seuils de widening}
\label{sec:point-fixe-composant-seuils}

Le widening des intervalles saute à $\pm\infty$ dès qu'une borne bouge ; avec
lui, le compteur gardé \code{if (count < 10) setCount(count + 1)} vaudrait
$[0, +\infty]$ alors que la garde le borne à 10. Le \emph{widening à seuils}
(\cref{def:widening}) saute à la place au plus petit seuil qui englobe la
nouvelle borne, et à $\pm\infty$ seulement au-delà du dernier seuil. Les
seuils sont récoltés une fois, après expansion :

\rustsnippet[label={lst:point-fixe-composant-thresholds}]{src/engine/fixpoint.rs}{1185}{1224}{la récolte des seuils : tous les littéraux numériques, y compris dans les fermetures}

\begin{decision}[ADR-014 --- widening à seuils, et pas de narrowing]
\emph{Contexte} (\adr{14}, 2026-06-27). \code{Interval::widen} jette toute
l'information numérique ; l'information qui bornerait la croissance (le
littéral \code{10} de la garde) est pourtant dans le programme.
\emph{Décision, partie 1.} Un ensemble fini de seuils $T$ par composant ;
\code{widen_to} ajouté à côté de \code{widen} sans changer la signature du
trait, en ne traitant que la composante numérique ; les seuils sont possédés
par le pilote du point fixe et passés en argument, jamais stockés dans les
valeurs. \emph{Partie 2, abandonnée} (révision du 2026-06-28). Le narrowing
classique --- une phase descendante après le post-point fixe --- n'a pas été
implémenté : le raffinement de branche pendant la montée plus le widening à
seuils, y compris sur l'arc arrière interne d'\code{analyze_cfg}, retrouvent
déjà toutes les bornes littérales ; et une descente ne pourrait de toute
façon pas faire redescendre le store d'état, dont les écritures s'accumulent
par join monotone. Trois conséquences annoncées ne sont donc pas réalisées :
\code{WidenOutcome}, \code{Config::narrow_iterations}, et la suppression du
recalcul \code{effect_setter_writes}, que l'ADR qualifie de \og crude
precision-recovery hack\fg{} et qui est resté
(\cref{sec:point-fixe-composant-pure}).
\end{decision}

Deux écarts entre le code et sa documentation sont à connaître. L'ADR
annonçait pour seuils les littéraux des \emph{gardes} et des inits ; le code
prend \emph{tous} les littéraux numériques du rendu, des inits de
\code{State}, des corps d'\code{Effect} et de \code{Handler}, en descendant
dans les \code{FnLit} (l.~1217--1219). C'est sans conséquence de soundness
(\og over-collecting is harmless\fg, l.~1189--1190) : l'ensemble reste fini,
un seuil de plus n'est qu'une borne candidate de plus. Inversement, la doc de
la fonction annonce \og all hook bodies\fg{} (l.~1187), mais les corps de
\code{Memo} et de \code{Callback} et les arguments des \code{Custom} ne sont
pas parcourus (le \code{match} des l.~1195--1201 les ignore). La
\cref{tab:point-fixe-composant-seuils} donne les seuils des exemples du
chapitre.

\begin{table}[htbp]\centering\small
\begin{tabular}{@{}lll@{}}
\toprule
composant & littéraux du programme & $T$ \\
\midrule
\code{Runaway} & \code{useState(0)}, \code{x + 1} & $\{0, 1\}$ \\
\code{Counter} & \code{useState(0)}, \code{count + 1} (dans le \code{onClick}) & $\{0, 1\}$ \\
\code{Guarded} & \code{useState(0)}, \code{count < 10}, \code{count + 1} & $\{0, 1, 10\}$ \\
\code{Lazy} & \code{useState(() => 7)}, \code{v < 20}, \code{v + 1} & $\{1, 7, 20\}$ \\
\code{Dead} & \code{useState(0)}, \code{const k = 3}, \code{k > 10}, \code{n + 1} & $\{0, 1, 3, 10\}$ \\
\code{Loader} & \code{useState(0)}, \code{setN(42)} & $\{0, 42\}$ \\
\bottomrule
\end{tabular}
\caption{Les seuils récoltés par \code{collect_thresholds} sur les exemples
du chapitre (relevés par la sonde). Le \code{7} de \code{Lazy} vient du
corps du \code{FnLit} initialiseur ; le \code{3} de \code{Dead} n'est pas une
garde mais une constante locale, et bornera pourtant la croissance d'un
palier (\cref{sec:point-fixe-composant-trous}).}
\label{tab:point-fixe-composant-seuils}
\end{table}

\subsection{Les corps de \code{useCallback}}

Un \code{useCallback} est réécrit par le lowering en \code{CallbackVal(l)} :
son corps quitte l'arbre d'expressions et vit dans l'entrée
\code{HookEntry::Callback}. Pour qu'un appel à travers la variable liée
(\code{const cb = useCallback(...); onLoad={(e) => cb(e)}}) exécute encore
le corps pour ses effets de bord, la préparation construit une table
\code{callback_bodies} des corps par label, exposée à l'interpréteur par
\code{QueryContext::callback_body} (\src{src/engine/fixpoint.rs}{282}{295},
\cref{chap:transfert-stores}). C'est le seul rôle du
\code{FixpointCtx} construit à chaque passe : ses deux autres champs,
\code{state} et \code{memo}, ne sont lus par aucune méthode du trait, et
semblent résiduels.

\subsection{L'amorçage des slots}
\label{sec:point-fixe-composant-amorcage}

Le sujet du point fixe est déclaré, avec ses compagnons, juste avant la
boucle :

\rustsnippet[label={lst:point-fixe-composant-porteurs}]{src/engine/fixpoint.rs}{297}{312}{le store sujet du point fixe et les accumulateurs de la boucle}

Puis chaque \code{useState} reçoit la valeur de son initialiseur :

\rustsnippet[label={lst:point-fixe-composant-seed}]{src/engine/fixpoint.rs}{314}{353}{l'amorçage : un join par \code{useState}, l'initialiseur paresseux exécuté}

L'init est évalué dans \code{initial_env} (constantes et props venues du
parent), contre des stores \emph{vides} et un tas neuf : un initialiseur ne
lit pas l'état, il le crée. Le cas particulier des l.~342--346 mérite un
rappel de React.

\begin{react}[Initialiseur paresseux]
\code{useState(() => e)} appelle le thunk \emph{une fois}, au montage, et
stocke sa valeur de retour ; aux rendus suivants l'argument est ignoré
(\og React saves the initial state once and ignores it on the next
renders\fg, \cite{ReactDocs}). L'état n'est jamais la fermeture elle-même.
Abstraire le \code{FnLit} comme une référence (ce que faisait une version
antérieure) faisait de chaque slot à init paresseux un \og dep instable\fg{}
--- un faux positif de corpus, cité par le commentaire. Le moteur exécute
donc le corps du thunk (\code{exec_body}) et amorce le slot avec son retour.
La règle \regle{lazy-init} exploite l'autre moitié du fait : un argument
\emph{non} paresseux est ré-évalué à chaque rendu pour rien.
\end{react}

\begin{exemple}{Un initialiseur paresseux et le seuil 20}{point-fixe-composant-lazy}
\begin{lstlisting}[language=TypeScript,caption={\file{/tmp/ch13/lazy.tsx}},label={lst:point-fixe-composant-lazy}]
import { useState, useEffect } from "react";

export function Lazy() {
  const [v, setV] = useState(() => 7);
  useEffect(() => { if (v < 20) setV(v + 1); }, [v]);
  return <div>{v}</div>;
}
\end{lstlisting}
L'IR de l'init est \texttt{FnLit \{ params: [], body\_cfg: Return(Lit(Int(7))) \}}.
L'amorçage donne \code{s0=num[7,7]}, pas une référence. Les tours donnent
$[7,8]$, $[7,9]$ (joins), puis au troisième changement le widening à seuils
cherche le plus petit seuil $\geq 10$ dans $T = \{1, 7, 20\}$ : c'est 20, d'où
$[7,20]$, stable au tour suivant. Sonde : \code{iterations = 3},
\code{state_store = s0=num[7,20]}, \code{effect_setter_writes = s0=num[8,20]},
\texttt{widen\_trace = [(0, WidenEvent \{ iteration: 3, writers: [1] \})]}. Le
CLI n'émet que l'Info \regle{widening-info} et conclut
\code{✓ 1 file(s) no issues found.} : le slot a été widené, mais la valeur
que l'effet écrit est bornée (\cref{sec:point-fixe-composant-pure}).
\end{exemple}

\begin{piege}
Le cas paresseux n'est reconnu que pour un \code{FnLit} \emph{sans paramètre
écrit en place}. Avec \code{const makeInit = () => 7;} au niveau du module
puis \code{useState(makeInit)}, l'init est un \code{Var} : il est évalué comme
une expression ordinaire, pas exécuté. Et \code{makeInit} n'est pas une
constante de module au sens du lowering. \code{ModuleConstInit} ne retient
que les littéraux primitifs, les objets, tableaux, \code{new}, regex et JSX,
et \code{createContext} (\src{src/lowering/mod.rs}{314}{336}), pas une
fonction fléchée. La lecture de \code{makeInit} manque donc l'environnement
et vaut $\Top$. Observé sur \file{/tmp/ch13/makeinit.tsx} (composant
\code{LazyVar}) avec le même effet \code{if (v < 20) setV(v + 1)} : la sonde
donne l'amorçage \code{s0=⊤}, \code{iterations = 0}, et le CLI émet un
\Warning{} \regle{infinite-loop} (\og may store a fresh reference into state
`v`\fg) : un faux positif. L'amorçage reste sound, car $\Top$ contient 7 ;
c'est seulement une perte de précision.
\end{piege}

\begin{invariant}{Tout slot est présent dès l'amorçage}{point-fixe-composant-amorcage}
Après l'amorçage, chaque label \code{State} du composant (hooks greffés
compris) est une clé du store, même si sa valeur est $\Bot$ (un init qui
s'évalue à $\Bot$ est rarissime, mais \code{update} insère la clé de toute
façon). Une passe qui part de \code{state.clone()} rend donc un store dont
\code{labels()} liste \emph{tous} les slots, qu'elle les ait écrits ou non.
Ce fait anodin explique une imprécision de provenance décrite au
\cref{sec:point-fixe-composant-widen-trace}.
\end{invariant}

% ===========================================================================
\section{Un tour de la boucle externe}
\label{sec:point-fixe-composant-tour}

La boucle (\src{src/engine/fixpoint.rs}{355}{530}) est une itération de
Kleene \emph{globale} (\cref{def:point-fixe}) : chaque tour ré-exécute tous
les corps, dans un ordre fixe, à partir du même store $S$. Il n'y a pas de
worklist à ce niveau --- la worklist n'existe qu'à l'intérieur d'un CFG
(\cref{chap:cfg-analyzer-dominance}). Chaque tour commence par
\code{let mut state_store = state.clone();} (l.~356) : c'est cette copie, et
non \code{state}, que toutes les passes du tour reçoivent et enrichissent.

\subsection{La passe de rendu}

\rustsnippet[label={lst:point-fixe-composant-render}]{src/engine/fixpoint.rs}{355}{380}{la passe de rendu d'un tour}

L'environnement d'entrée est \code{initial_env} : les constantes de module et,
en analyse inter, les props que le parent a liées. \code{analyze_cfg} rend les
environnements de sortie de chaque bloc (\code{bs}, rangé dans
\code{block_states} malgré son nom, \cref{chap:cfg-analyzer-dominance})
et un store \code{state_from_render}, noté $R$, qui part de
\code{state_store} et y joint chaque écriture rencontrée pendant le rendu.
Comme \code{analyze_cfg} part d'une copie du store reçu et n'y fait que
joindre (\srcline{src/engine/cfg_analyzer.rs}{49},
\cref{chap:cfg-analyzer-dominance}), $R \lgeq S$ toujours.
Deux sortes d'écritures peuvent survenir pendant un rendu : un setter appelé
dans le corps (\regle{setter-in-render} s'en occupe), et, quand \code{inter}
est présent, l'analyse d'un enfant \code{<Child .../>} rencontré dans le JSX,
lancée \emph{pendant} le rendu par \code{eval_comp_app}
(\cref{chap:transfert-stores}) ; l'enfant peut alors écrire dans le
\code{SharedStateStore}, que le test de convergence importera.

\subsection{Les mémos : une valeur qui n'est qu'une stabilité}
\label{sec:point-fixe-composant-memos}

\rustsnippet[label={lst:point-fixe-composant-memo}]{src/engine/fixpoint.rs}{382}{413}{le recalcul des mémos depuis l'environnement de sortie du rendu}

L'environnement de sortie du rendu, \code{env_exit}, est le join des
environnements de sortie des blocs terminés par \code{Return} :

\rustsnippet[label={lst:point-fixe-composant-exitenv}]{src/engine/fixpoint.rs}{1226}{1237}{\code{exit_env} : le join des sorties, par \code{reduce} et non par \code{fold}}

Le \code{reduce} plutôt que \code{fold(bottom, join)} n'est pas une
coquetterie : \code{AbstractEnv::bottom()} n'est pas neutre pour le join
(une variable présente d'un seul côté devient $\Top$,
\cref{chap:cfg-analyzer-dominance}) ; un CFG sans \code{Return}
atteint rend un environnement vide, dont chaque lecture vaut $\Top$. C'est
sound, et cette fonction est dupliquée à l'identique dans
\code{AnalysisResult::exit_env} (\src{src/engine/analysis_result.rs}{274}{289}).

Le recalcul lui-même est délégué au domaine
(\code{Transfer::recompute_memo}, \cref{chap:transfert-stores}) : la
valeur d'un \code{useMemo} ou d'un \code{useCallback} est
\emph{uniquement une stabilité de référence} --- \code{Stable} pour des
deps \code{[]} connues vides, \code{Unknown} pour des deps illisibles,
sinon le join des stabilités de ses deps, un dep \code{StateVal(l)} donnant
\code{Versioned({(composant, l)})}. Le corps du mémo n'est \emph{jamais}
exécuté par le moteur : \code{useMemo(() => x + 1, [x])} n'a pas de
composante numérique. Deux détails de la boucle en découlent.

\begin{itemize}
  \item Les écritures dans le memo store sont \emph{différées} (commentaire des
    l.~383--387) : chaque \code{recompute_memo} lit un instantané cohérent,
    et une chaîne mémo $\to$ mémo converge \og à travers les itérations\fg,
    avec un tour de retard par maillon.
  \item Le rendu du tour $n$ lit les mémos calculés au tour $n-1$, puisque le
    recalcul suit la passe de rendu. Au tour 0, il n'y a encore rien, et
    \code{MemoStore::get} d'un label absent vaut $\Top$ :
\end{itemize}

\rustsnippet[label={lst:point-fixe-composant-memoget}]{src/domains/stores/memo_store.rs}{27}{30}{un mémo pas encore calculé vaut $\Top$ --- à l'inverse du \code{StateStore}, où l'absent vaut $\Bot$}

\begin{exemple}{Une boucle portée par un \code{useMemo}}{point-fixe-composant-memoloop}
\begin{lstlisting}[language=TypeScript,caption={\file{/tmp/ch13/memoloop.tsx}},label={lst:point-fixe-composant-memoloop}]
import { useState, useEffect, useMemo } from "react";

export function MemoLoop() {
  const [x, setX] = useState(0);
  const d = useMemo(() => x + 1, [x]);
  useEffect(() => {
    setX(d);
  }, [d]);
  return <div>{x}</div>;
}
\end{lstlisting}
Concrètement le composant boucle ($0 \to 1 \to 2 \to \dots$). Abstraitement,
au tour 0 le rendu lie \code{d} à \code{MemoVal(1)} avant tout recalcul, soit
$\Top$ ; l'effet exécute \code{setX(⊤)} et le slot saute à $\Top$ en un
tour ; le tour 1 ne change rien. La sonde donne \code{iterations = 1},
\code{state_store = s0=⊤}, \code{widen_trace = []}, et après convergence le
mémo vaut \texttt{ref:Versioned(\{(C0, 0)\})}, d'où
\code{effect_setter_writes = s0=ref:Versioned(...)} et un trigger
\code{exact: false}. Le CLI répond \code{✓ 1 file(s) no issues found.} : un
faux négatif, que nous rangerons au \cref{sec:point-fixe-composant-trous}
(\issue{157}). La leçon pour la boucle : \emph{un slot qui atteint $\Top$ par
join n'a pas été widené}, donc n'est pas dans \code{widen_trace}, donc n'est
pas un signal de divergence.
\end{exemple}

\begin{piege}
Le memo store n'entre pas dans le test de convergence : seul le
\code{StateStore} y est comparé. Le memo store est recalculé à chaque tour et
son dernier contenu est publié, mais si un mémo change au dernier tour sans
que l'état change, aucun tour supplémentaire ne rejoue les effets avec sa
nouvelle valeur. Les effets du tour ne la voient pas non plus : ils lisent
la variable liée par le rendu, donc le mémo du recalcul précédent. Sur
\code{MemoLoop}, l'effet du tour 0 écrit bien $\Top$, pas la valeur
\code{Versioned} calculée juste avant lui. La question est donc de savoir si
une valeur de mémo peut \emph{monter} entre deux tours sans que l'état
bouge. Les deux sources de variation connues descendent. $\Top$ au tour 0
laisse place à une stabilité plus précise. Dans une chaîne \code{m2}
$\to$ \code{m1}, \code{m2} garde un tour de retard, mais sur une valeur de
\code{m1} qui elle-même ne fait que descendre. Dans ces deux cas, la valeur
périmée qu'une écriture a jointe dans l'état est au-dessus de la valeur
finale, et l'état reste sound. Un dep dérivé de l'état reste la seule voie de
montée. Un dep qui \emph{est} un slot n'en fait pas partie, car il vaut
toujours \code{Versioned} (\src{src/domains/transfer/state_value.rs}{79}{88}).
Au commit \snapshotCommit, nous n'avons pas construit d'instance d'une telle
montée. La passe de rafraîchissement du \cref{sec:point-fixe-composant-post}
corrige de toute façon ce que les \emph{règles} lisent.
\end{piege}

\subsection{Les effets : tous, à chaque tour, depuis la sortie du rendu}
\label{sec:point-fixe-composant-effets}

\rustsnippet[label={lst:point-fixe-composant-effets}]{src/engine/fixpoint.rs}{415}{449}{les passes d'effets d'un tour}

Trois décisions tiennent dans ces trente lignes.

\paragraph{Tous les effets, quelles que soient leurs deps.} Le motif
\texttt{HookEntry::Effect \{ label, body\_cfg, .. \}} ignore \code{deps}. Un effet
\code{[]} de montage, un effet \code{[count]}, un effet sans tableau sont
exécutés de la même façon, à chaque tour. C'est une sur-approximation de la
règle de React \og l'effet tourne si et seulement si un dep a changé selon
\code{Object.is}\fg{} (\cref{chap:react-modele}) : l'ensemble des écritures
possibles est un sur-ensemble de ce que React ferait. La précision est
récupérée \emph{dans les règles}, qui relisent les deps : \regle{infinite-loop}
saute les effets à \code{[]} et ceux dont tous les deps sont prouvés stables
(\src{src/rules/impls/infinite_loop.rs}{96}{113}), et la relation
\relation{effect_triggers} dit lesquels un slot relance. L'\adr{4} décrivait
au contraire \og for each Effect whose decision = Effect\fg{} ; le code a
choisi plus large, et il n'a pas de notion de décision.

\paragraph{Depuis \code{env_exit}.} Chaque effet part de l'environnement de
sortie du rendu de \emph{ce} tour : il voit les valeurs que ce rendu a
liées, exactement comme la fermeture que React a créée pendant le rendu.

\paragraph{Depuis le même \code{state_store}.} Tous les effets du tour
reçoivent la même copie \code{state_store} ; ce qu'un effet écrit n'est pas
visible des effets suivants du même tour, seulement du tour suivant, via
le join \code{state_from_effects}. L'ordre concret des effets (enfants
d'abord, puis ordre syntaxique) est abstrait par un join commutatif.

\begin{react}[Fermetures, file de mises à jour et batching]
Un effet est une fermeture créée pendant le rendu : il lit les variables du
rendu qui l'a créé, pas les valeurs courantes (\textsc{Eff} de React-tRace
met en file $\langle \lambda\_.e, \sigma\rangle$ avec l'environnement
$\sigma$ \emph{du rendu}, \cref{chap:react-modele}). Un \code{setX(v)}
n'écrit pas l'état : il met \code{v} en file, et React applique la file au
rendu suivant --- depuis React 18, toutes les mises à jour d'un même effet ou
d'un même événement sont regroupées (\emph{batching}). Le moteur abstrait ce
double mécanisme en un seul : la mise à jour faible
$S[\ell] := S[\ell] \join v$, qui contient tout résultat possible de la file
(dernier gagne, updaters enchaînés, bail-out \code{Object.is}), au prix de
l'ordre. Un effet qui appelle \code{setM(doubled)} puis
\code{setM(prev => prev + 1)} voit, dans l'updater, un \code{prev} qui
contient déjà la première écriture de la passe --- c'est la file, abstraite.
\end{react}

\begin{exemple}{Un effet de montage, rejoué à chaque tour}{point-fixe-composant-loader}
\begin{lstlisting}[language=TypeScript,caption={\file{/tmp/ch13/loader.tsx}},label={lst:point-fixe-composant-loader}]
import { useState, useEffect } from "react";

export function Loader() {
  const [n, setN] = useState(0);
  useEffect(() => {
    setN(42);
  }, []);
  return <div>{n}</div>;
}
\end{lstlisting}
React exécute l'effet une fois. Le moteur l'exécute à chaque tour :
\begin{center}\small
\begin{tabular}{@{}lllll@{}}
\toprule
tour & $R$ & effet 1 & $N$ & action \\
\midrule
amorçage & & & & \code{s0=num[0,0]} \\
0 & $[0,0]$ & $[0,42]$ & $[0,42]$ & \code{iteration = 1} $< 3$ : $S := N$ \\
1 & $[0,42]$ & $[0,42]$ & $[0,42]$ & $N \lleq S$ : sortie \\
\bottomrule
\end{tabular}
\end{center}
Résultat : \code{iterations = 1}, \code{s0=num[0,42]} --- l'enveloppe
convexe de $\{0, 42\}$, que le domaine des intervalles ne peut pas
distinguer de $\{0, \dots, 42\}$ ---, \code{effect_setter_writes =
s0=num[42,42]}, aucun trigger (deps \code{[]}). La ré-exécution au tour 1
est sans conséquence ici (l'écriture est idempotente). Le CLI émet un
\Warning{} \regle{unnecessary-rerender} (\og mount-only effect sets state
`n` to a constant different from its initial value\fg) et
\code{verified infinite-loop}.
\end{exemple}

\begin{exemple}{Deux effets qui se nourrissent : le batching abstrait}{point-fixe-composant-mutual}
\begin{lstlisting}[language=TypeScript,caption={\file{/tmp/ch13/mutual.tsx}},label={lst:point-fixe-composant-mutual}]
import { useState, useEffect } from "react";

export function Mutual() {
  const [x, setX] = useState(0);
  const [y, setY] = useState(0);
  useEffect(() => { setY(x + 1); }, [x]);
  useEffect(() => { setX(y + 1); }, [y]);
  return <div>{x + y}</div>;
}
\end{lstlisting}
Au tour 0, les deux effets partent du même store $\{s_0 = [0,0], s_1 = [0,0]\}$ :
\begin{sortie}
    iter 0: effect 2 out = s0=num[0,0], s1=num[0,1]
    iter 0: effect 3 out = s0=num[0,1], s1=num[0,0]
    iter 0: render out = s0=num[0,0], s1=num[0,0]
    iter 0: new_state = s0=num[0,1], s1=num[0,1]
\end{sortie}
L'effet 3 ne voit pas l'écriture de l'effet 2 : il lit $y = [0,0]$ et écrit
$x \in [0,1]$. Chaque tour n'ajoute donc que $+1$ à chaque slot, et le
troisième changement déclenche le widening des deux :
\code{s0=num[0,inf], s1=num[0,inf]}, \code{widen_trace} sur les labels 0 et
1. Le CLI émet deux \Warning{} \regle{infinite-loop} (un par slot) et deux
\regle{derived-state}.
\end{exemple}

\subsection{Les handlers : dans le point fixe, hors du signal}
\label{sec:point-fixe-composant-handlers}

\rustsnippet[label={lst:point-fixe-composant-handlers}]{src/engine/fixpoint.rs}{451}{483}{les passes de handlers : même schéma que les effets, avec le commentaire décisif}

Le lowering a extrait chaque \code{onClick={...}} et chaque
\code{addEventListener(ev, fn)} d'un effet en une entrée
\code{HookEntry::Handler} avec son propre CFG (\cref{chap:lowering-cfg}).
Les handlers sont analysés exactement comme les effets --- même
\code{env_exit}, même \code{state_store} --- et leur store $H$ est joint
dans $N$. Mais leur croissance n'entre pas dans \code{widen_trace}. Le
commentaire des l.~452--453 dit tout : \og Handlers run 0..N times $\to$
include in fixpoint for sound range approx. NOT tracked in widened\_labels
(handler-caused widening $\neq$ InfiniteLoop)\fg.

\begin{react}[Événements]
Un handler s'exécute sur une entrée extérieure, zéro, une ou $N$ fois, dans
un ordre arbitraire par rapport aux autres handlers ; en React-tRace c'est
\textsc{StepEvent}, qui n'est franchi qu'en mode \emph{event loop}
(\cref{chap:react-modele}). Un compteur que l'utilisateur
incrémente mille fois atteint mille : c'est une \emph{valeur} que l'analyse
doit admettre, mais ce n'est pas une boucle, puisque chaque tour exige un
clic. Un effet qui incrémente, lui, relance le cycle sans intervention.
\end{react}

\begin{decision}[ADR-009 --- points d'entrée et classes de déclenchement]
\emph{Décision} (\adr{9}). Un composant est un ensemble de \emph{points
d'entrée} --- rendu, effets, callbacks différés (\code{.then}, timers),
handlers --- analysés par la même machinerie, et qui ne diffèrent que sur deux
axes : \og dans le cycle automatique ?\fg{} et \og un widening induit est-il
un bug ?\fg. Le rendu \emph{est} le cycle ; un effet ou un \code{.then} y sont
(oui / oui) ; un handler n'y est pas (non / non : \og clicking 1000$\times$
isn't a bug\fg). La §5 de l'ADR (\emph{multiplicity}) demande que les
handlers soient tout de même dans la boucle de point fixe, pour la
soundness des intervalles, avec widening, mais exclus du signal ---
c'est le code des l.~451--483 et le filtre \code{new_state_incycle} du test
de convergence. \emph{Alternative refusée} : descendre uniformément dans
tout callback, qui produirait un faux positif \regle{infinite-loop} sur
\code{addEventListener('click', () => setCount(c => c + 1))}. \emph{Un choix
en tension} : pour un callee \code{Unknown}, l'ADR retient \code{skip}
(\og we accept the FN\fg), ce que l'invariant du projet interdit en principe ;
le mécanisme B6 du \cref{chap:transfert-stores} en réduit la portée.
\end{decision}

\begin{exemple}{Un compteur par clic : widené, pas tracé}{point-fixe-composant-counter}
\begin{lstlisting}[language=TypeScript,caption={\file{/tmp/ch13/counter.tsx}},label={lst:point-fixe-composant-counter}]
import { useState } from "react";

export function Counter() {
  const [count, setCount] = useState(0);
  return <button onClick={() => setCount(count + 1)}>{count}</button>;
}
\end{lstlisting}
IR : \code{hook 0: State init=Lit(Int(0))}, \code{hook 1: Handler click}.
Seuils $\{0, 1\}$.
\begin{center}\small
\begin{tabular}{@{}llllp{52mm}@{}}
\toprule
tour & $R$ & handler 1 & $N$ & action \\
\midrule
0 & $[0,0]$ & $[0,1]$ & $[0,1]$ & \code{iteration = 1} : $S := N$ \\
1 & $[0,1]$ & $[0,2]$ & $[0,2]$ & \code{iteration = 2} : $S := N$ \\
2 & $[0,2]$ & $[0,3]$ & $[0,3]$ & \code{iteration = 3} : \code{widen_to} $\to [0,+\infty]$ ; labels changés dans $R \join E$ : aucun \\
3 & $[0,+\infty]$ & $[0,+\infty]$ & $[0,+\infty]$ & $N \lleq S$ : sortie \\
\bottomrule
\end{tabular}
\end{center}
Sonde : \code{iterations = 3}, \code{state_store = s0=num[0,inf]},
\code{widen_trace = []}. CLI : \code{Counter: 3 iteration(s), widened: []},
aucun diagnostic. Le slot est bien widené (l'intervalle $[0,+\infty]$ est
la seule valeur sound pour un compteur de clics), mais comme seul le handler
l'a fait bouger, \code{changed_labels} de $R \join E$ est vide et rien
n'entre dans \code{widen_trace}.
\end{exemple}

L'inverse est tout aussi important : les handlers font partie du point fixe
\emph{pour rendre les effets sound}. Le test unitaire
\code{handler_enables_infinite_loop_detection}
(\src{src/engine/fixpoint.rs}{2564}{2694}) analyse un effet
\code{if (count > 1) setCount(count + 1)} avec deps \code{[count]}, plus un
handler \code{setCount(count + 1)}. Sans les handlers dans la boucle, le
slot resterait à $[0,0]$. Dans la branche, \code{count} serait raffiné à
$\Bot$ (\code{narrow_gt(1)} sur $[0,0]$), l'écriture \code{count + 1} vaudrait
$\Bot$, et la boucle infinie que
déclenche le \emph{deuxième} clic ne serait jamais vue : un faux négatif.
Avec eux, le handler fait croître \code{count} tour après tour jusqu'à
rendre la branche vivante ; l'effet écrit alors, et c'est \emph{lui} qui
fait entrer le label dans \code{widen_trace}. Les handlers participent à la
\emph{valeur} sans participer au \emph{signal}.

\subsection{Le join et le test de convergence}
\label{sec:point-fixe-composant-join}

\rustsnippet[label={lst:point-fixe-composant-test}]{src/engine/fixpoint.rs}{485}{497}{le join des quatre sources et le test d'arrêt}

Le nouvel état est le join de quatre stores : $R$ (rendu), $E$ (effets),
$H$ (handlers) et la tranche du \code{SharedStateStore} qui concerne ce
composant --- les écritures faites, pendant ce tour ou un tour antérieur,
par un setter de ce composant appelé \emph{ailleurs} (un enfant qui a reçu
\code{setCount} en prop). C'est ici, et non en tête de boucle, que
l'étape 1 du plan de \cref{lst:point-fixe-composant-signature} a lieu.
Le store partagé n'a qu'une opération d'extraction :

\rustsnippet[label={lst:point-fixe-composant-slice}]{src/domains/stores/shared_state_store.rs}{57}{66}{la tranche d'un composant dans le \code{SharedStateStore}}

et le test d'arrêt est l'ordre point à point des stores :

\rustsnippet[label={lst:point-fixe-composant-leq}]{src/domains/stores/state_store.rs}{69}{77}{$S \lleq S'$ : chaque label de $S$ est en dessous de sa valeur dans $S'$}

\code{leq_pointwise} s'appuie sur le \code{PartialOrd} de \code{StateValue}
(\code{Less | Equal}) : deux valeurs incomparables font échouer le test, ce
qui est la bonne réponse (le nouvel état n'est pas contenu dans l'ancien).

\begin{proposition}{Le test d'arrêt est un test d'égalité}{point-fixe-composant-egalite}
À chaque tour, $R \lgeq S$, $E \lgeq S$ et $H \lgeq S$, car chacune de ces
passes part de \code{state.clone()} et n'y fait que joindre
(\srcline{src/engine/cfg_analyzer.rs}{49}) ; donc $N \lgeq S$. Par
antisymétrie, \code{new_state.leq(&state)} équivaut à $N = S$ sur tous les
labels : la boucle s'arrête exactement quand un tour de plus n'a rien
ajouté. La suite $(S_k)$ est croissante, et $S_k$ à l'arrêt est un
post-point fixe de la fonction \og un tour\fg{} au sens de la
\cref{def:point-fixe-composant-etat}.
\end{proposition}

La tranche partagée réserve une surprise que le \cref{sec:point-fixe-composant-trous}
détaillera : quand \code{inter} est présent, un composant écrit \emph{ses
propres} slots dans le \code{SharedStateStore}. Sur \code{Counter} analysé
par \code{analyze_program}, la sonde lit
\texttt{shared\_state = \{(ComponentId(0), 0): number[1, inf]\}} --- l'écriture
du handler, rangée dans le store partagé par le second bras de
\code{exec_setter_call} (\src{src/domains/interp/interpreter.rs}{368}{391}),
puis réimportée par \code{slice}. Ici c'est sans effet (la valeur est déjà
dans $H$) ; avec un updater fonctionnel, ce ne l'est plus.

% ===========================================================================
\section{Convergence, widening et terminaison}
\label{sec:point-fixe-composant-convergence}

\subsection{Le compteur et les deux opérateurs}

Quand $N \not\lleq S$, la boucle doit choisir comment avancer :

\rustsnippet[label={lst:point-fixe-composant-widen}]{src/engine/fixpoint.rs}{499}{530}{join tant que le compteur est bas, widening à seuils ensuite, garde-fou à cent}

Lecture pas à pas.

\begin{enumerate}
  \item \emph{Le compteur} \code{iteration} n'est incrémenté que si l'état a
    changé. Au $k$-ième tour (compté de 0) qui change l'état, il passe à
    $k+1$. Avec \code{widen_threshold = 3}, les deux premiers changements
    sont des joins ($S := N$), le troisième et les suivants des widenings.
    Le compteur est \emph{global au composant}, pas par label : un slot qui
    commence à croître tard est widené dès son premier changement si le
    seuil est déjà atteint.
  \item \emph{Le widening} \code{widen_to(&new_state, &thresholds)} porte
    sur \emph{tout} le store, y compris la croissance venue des handlers ;
    mais \code{widen_trace} ne retient que les labels changés dans
    \code{new_state_incycle} $= R \join E$ (l.~518). C'est la traduction
    exacte de l'\adr{9} : la valeur est sound, le signal ne compte que le
    cycle.
  \item \emph{La première itération gagne} : \code{or_insert_with} garde
    l'événement du premier widening d'un label, le plus informatif pour la
    chaîne de témoins (\adr{19}).
  \item \emph{Le garde-fou} : à \code{iteration >= 100}, un widening
    \emph{sans seuils} (\code{widen}, pas \code{widen_to}) sur tout le store,
    tous les labels du store dans \code{widen_trace} (y compris ceux qui
    n'ont pas bougé), et \code{break} immédiat.
\end{enumerate}

Au niveau du store, les trois opérateurs ne diffèrent que par la fonction
appliquée à chaque label :

\rustsnippet[label={lst:point-fixe-composant-merge}]{src/domains/stores/state_store.rs}{35}{62}{\code{join}, \code{widen}, \code{widen_to} : un seul parcours, trois combinateurs}

et au niveau de la valeur, seule la composante numérique a des seuils :

\rustsnippet[label={lst:point-fixe-composant-svwiden}]{src/domains/impls/state_value.rs}{657}{676}{le widening du produit : intervalles et chaînes ont un vrai widening, le reste est un join}

L'\cref{alg:point-fixe-composant-boucle} récapitule la boucle complète, garde-fou
et provenance compris ; c'est la version détaillée de
l'algorithme réduit du \cref{chap:react-modele}.

\begin{algorithm}[htbp]
\Entree{\code{render_cfg}, \code{hooks} (après expansion), \code{initial_env}, seuils $T$, $t$ = \code{widen_threshold}}
\Sortie{$S$ convergé, \code{iteration}, \code{widen_trace}}
$S \leftarrow \Bot$ ; \PourCh{$\mathtt{State}\{\ell, \mathit{init}\}$}{$S[\ell] \leftarrow S[\ell] \join \sem{\mathit{init}}$}
$i \leftarrow 0$ ; $\mathit{trace} \leftarrow \emptyset$\;
\Tq{vrai}{
  $S' \leftarrow$ copie de $S$\;
  $(\mathit{bs}, R) \leftarrow \mathtt{analyze\_cfg}(\mathtt{render\_cfg}, \mathtt{initial\_env}, S')$ ; $\mathit{env} \leftarrow \mathtt{exit\_env}(\mathit{bs})$\;
  \PourCh{$\mathtt{Memo}/\mathtt{Callback}\{\ell, \mathit{deps}\}$}{$\mathit{memo}[\ell] \leftarrow \mathtt{recompute\_memo}(\mathit{deps}, \mathit{env})$ \tcp*{écritures différées}}
  $E \leftarrow \Bot$ ; $W \leftarrow \emptyset$\;
  \PourCh{$\mathtt{Effect}\{e, \mathit{body}\}$}{$(\_, E_e) \leftarrow \mathtt{analyze\_cfg}(\mathit{body}, \mathit{env}, S')$ ; $E \leftarrow E \join E_e$ ; $W[\ell] \mathrel{+}= e$ pour tout $\ell \in \mathtt{labels}(E_e)$\;}
  $H \leftarrow \Bot$ ; \PourCh{$\mathtt{Handler}\{h, \mathit{body}\}$}{$(\_, H_h) \leftarrow \mathtt{analyze\_cfg}(\mathit{body}, \mathit{env}, S')$ ; $H \leftarrow H \join H_h$\;}
  $N_{\mathrm{cycle}} \leftarrow R \join E$ ; $N \leftarrow N_{\mathrm{cycle}} \join H \join \mathtt{shared.slice}(\mathit{comp})$\;
  \Si{$N \lleq S$}{\textbf{sortir}\;}
  $i \leftarrow i + 1$\;
  \Si{$i \geq 100$}{$\mathit{trace}[\ell] \leftarrow (i, W[\ell])$ pour tout $\ell \in \mathtt{labels}(S)$ non déjà tracé ; $S \leftarrow S \widen N$ ; \textbf{sortir}\;}
  \eSi{$i \geq t$}{
    \PourCh{$\ell \in \mathtt{changed\_labels}(N_{\mathrm{cycle}}, S)$ non déjà tracé}{$\mathit{trace}[\ell] \leftarrow (i, W[\ell])$\;}
    $S \leftarrow S \widen_T N$\;
  }{$S \leftarrow N$\;}
}
\caption{La boucle externe de \code{analyze_component_impl}
(\src{src/engine/fixpoint.rs}{297}{530}). $W$ est la table locale
\code{slot_writers} (\og quels effets ont un store contenant $\ell$\fg) ;
$\widen_T$ est \code{widen_to} avec les seuils $T$.}
\label{alg:point-fixe-composant-boucle}
\end{algorithm}

\subsection{Trois déroulés}
\label{sec:point-fixe-composant-deroules}

Les tableaux qui suivent ont été relevés avec la réplique de la sonde ; les
valeurs finales coïncident avec celles du moteur.

\paragraph{\code{Runaway} (\cref{lst:point-fixe-composant-runaway}).} Seuils
$\{0, 1\}$. L'effet lit \code{x} dans \code{env_exit} (la valeur de fin de
rendu) et écrit \code{x + 1} :

\begin{table}[htbp]\centering\small
\begin{tabular}{@{}lllp{54mm}@{}}
\toprule
tour & \code{env_exit(x)} & effet 1 & $N$ et action \\
\midrule
0 & $[0,0]$ & $[0,0] \join ([0,0] + 1) = [0,1]$ & $[0,1]$ ; \code{iteration = 1} : $S := [0,1]$ \\
1 & $[0,1]$ & $[0,1] \join [1,2] = [0,2]$ & $[0,2]$ ; \code{iteration = 2} : $S := [0,2]$ \\
2 & $[0,2]$ & $[0,3]$ & $[0,3]$ ; \code{iteration = 3} : \texttt{widen\_to([0,2],} \texttt{[0,3], \{0,1\})} : la borne haute 3 dépasse 2 et aucun seuil n'est $\geq 3$ : $+\infty$ ; \texttt{widen\_trace[0]} $=$ \texttt{\{3, [1]\}} \\
3 & $[0,+\infty]$ & $[0,+\infty]$ & $N \lleq S$ : sortie, \code{iterations = 3} \\
\bottomrule
\end{tabular}
\caption{Les tours de \code{Runaway}.}
\label{tab:point-fixe-composant-runaway}
\end{table}

\paragraph{\code{Guarded}.} Le compteur gardé, cas d'école de l'\adr{14} :

\begin{lstlisting}[language=TypeScript,caption={\file{/tmp/ch13/guarded.tsx}},label={lst:point-fixe-composant-guarded}]
import { useState, useEffect } from "react";

export function Guarded() {
  const [count, setCount] = useState(0);
  useEffect(() => {
    if (count < 10) setCount(count + 1);
  }, [count]);
  return <div>{count}</div>;
}
\end{lstlisting}

Seuils $\{0, 1, 10\}$. L'effet a quatre blocs (\code{Branch count < 10},
la branche qui écrit, la branche vide, la jonction) et le raffinement de
garde (\code{narrow_env_for_branch}, \cref{chap:cfg-analyzer-dominance}) borne \code{count} dans la
branche \emph{then} :

\begin{table}[htbp]\centering\small
\begin{tabular}{@{}llllp{58mm}@{}}
\toprule
tour & \code{count} & \code{count} dans le \emph{then} & effet & $N$ et action \\
\midrule
0 & $[0,0]$ & $[0,0]$ & $[0,1]$ & join \\
1 & $[0,1]$ & $[0,1]$ & $[0,2]$ & join \\
2 & $[0,2]$ & $[0,2]$ & $[0,3]$ & \code{iteration = 3} : plus petit seuil $\geq 3$ : \textbf{10} ; $S := [0,10]$ ; \code{widen_trace[0]} \\
3 & $[0,10]$ & \code{narrow_lt(10)} : $[0,9]$ & $[0,10] \join [1,10]$ & $[0,10] \lleq S$ : sortie \\
\bottomrule
\end{tabular}
\caption{Les tours de \code{Guarded} : le seuil 10 arrête la croissance, et la
garde raffinée rend l'écriture bornée.}
\label{tab:point-fixe-composant-guarded}
\end{table}

Résultat : \code{s0=num[0,10]}, \code{effect_setter_writes = s0=num[1,10]},
un trigger \texttt{\{hook: 1, dep: 0, slot: (C0, 0), exact: true\}}. Le CLI
n'émet que l'Info \regle{widening-info} et \code{verified infinite-loop} :
le label est bien dans \code{widen_trace}, mais l'écriture rejouée est
bornée (\cref{sec:point-fixe-composant-pure}).

\paragraph{\code{Dead}.} Un dernier déroulé montre ce qu'un seuil qui n'est
pas une garde fait à la suite :

\begin{lstlisting}[language=TypeScript,caption={\file{/tmp/ch13/dead.tsx}},label={lst:point-fixe-composant-dead}]
import { useState, useEffect } from "react";

export function Dead() {
  const [n, setN] = useState(0);
  const k = 3;
  useEffect(() => {
    if (k > 10) { setN(n + 1); }
  });
  return <div>{n}</div>;
}
\end{lstlisting}

\begin{sloppypar}
Concrètement, \code{k > 10} est toujours faux : l'effet n'écrit jamais.
Abstraitement, \code{narrow_env_for_branch} raffine \code{k} à $\Bot$ dans la
branche, mais le bloc est tout de même exécuté
(\cref{chap:cfg-analyzer-dominance}) ; \code{setN(n + 1)} lit \code{n},
qui n'est pas raffiné, et écrit. Seuils $\{0, 1, 3, 10\}$ : la suite est
$[0,1]$, $[0,2]$ (joins), puis $[0,3]$ au seuil 3, puis $[0,10]$ au seuil 10,
puis $[0,+\infty]$ ; \code{iterations = 5}, et \texttt{widen\_trace[0]} vaut
\texttt{\{iteration: 3, writers: [1]\}} (la première itération de widening,
pas la dernière). Le CLI émet un \Warning{} \regle{infinite-loop} : un faux
positif, toléré par l'invariant, dont la cause est au
\cref{chap:cfg-analyzer-dominance} et pas dans la boucle externe.
\end{sloppypar}

\subsection{\code{iterations} et \code{widen_trace}}
\label{sec:point-fixe-composant-widen-trace}

Deux champs de l'\code{AnalysisResult} racontent la convergence.
\code{iterations} est la valeur du compteur à la sortie : le nombre de tours
qui ont \emph{changé} l'état, pas le nombre de tours exécutés (il y en a
toujours un de plus). Un composant dont l'état ne bouge pas au premier tour
a \code{iterations = 0} : c'est le cas de \code{Fresh}
(\cref{lst:point-fixe-composant-fresh}), dont le store est à
\code{PerRender} dès l'amorçage. \code{widen_trace} est la table des labels
widenés par le cycle :

\rustsnippet[label={lst:point-fixe-composant-widenevent}]{src/engine/analysis_result.rs}{18}{28}{la provenance d'un widening}

Elle a deux lecteurs : \regle{widening-info}
(\srcline{src/rules/impls/widening_info.rs}{20}), qui émet une Info par
label, et la première porte de \regle{infinite-loop}
(\cref{sec:point-fixe-composant-pure}). Le champ \code{iteration} est exact.
Le champ \code{writers} ne l'est pas.

\begin{piege}
\code{writers} est censé lister \og les effets dont la passe a écrit ce slot
pendant l'itération du widening\fg. Il liste en réalité \emph{tous les
effets du composant}, pour chaque slot. La table locale \code{slot_writers}
(\cref{lst:point-fixe-composant-effets}, l.~444--446) est remplie pour chaque
label de \code{eff_state.labels()} ; or \code{eff_state} part de
\code{state_store.clone()}, qui contient tous les slots amorcés
(\cref{inv:point-fixe-composant-amorcage}). Sur \code{Mutual}
(\cref{lst:point-fixe-composant-mutual}), la sonde donne
\code{writers: [2, 3]} pour les deux slots, et la chaîne \code{--trace} du
CLI attribue à l'effet 3 une écriture de \code{y} qu'il ne fait pas :
\begin{sortie}
    warn   infinite-loop  [hook:1]  (line 6:2)  this effect keeps pushing state `y` (its deps do not provably gate it, so the effect can re-run every render) to new values on every run. Potential infinite render loop
       → state `y` is written here [hook:2] (line 6:2)
       → state `y` is written here [hook:3] (line 7:2)
       → the abstract value of state `y` kept growing and was widened at iteration 3
\end{sortie}
Aucun verdict n'en dépend : c'est une imprécision de \emph{provenance}, pas
de soundness. Le correctif naturel serait de ne retenir que
\code{eff_state.changed_labels(&state_store)}.
\end{piege}

\subsection{Terminaison}
\label{sec:point-fixe-composant-terminaison}

\begin{theoreme}{Terminaison de la boucle externe}{point-fixe-composant-terminaison}
Pour tout composant et tout ensemble fini de seuils $T$, la boucle de
\code{analyze_component_impl} s'arrête. Preuve. La suite $(S_k)$ est
croissante (\cref{prop:point-fixe-composant-egalite}) et, à partir du tour
$t$, chaque pas est un widening $S \widen_T N$ point à point sur des valeurs
\code{StateValue}. Il suffit donc que toute chaîne
$v_0 \lleq v_0 \widen_T v_1 \lleq \dots$ de \code{StateValue} se stabilise,
composante par composante (\cref{lst:point-fixe-composant-svwiden}) :
\begin{itemize}
  \item \code{num} : \code{Interval::widen_to} ne fait que relâcher les
    bornes vers un seuil de $T$ ou vers $\pm\infty$ ; chaque borne visite
    chaque seuil au plus une fois avant $\pm\infty$
    (\src{src/domains/impls/interval.rs}{112}{146}, \adr{14}) : au plus
    $2(|T| + 1)$ pas ;
  \item \code{str} : \code{StrConst::widen} bascule à $\Top$ au-delà de
    \code{STR_WIDEN_THRESHOLD = 4} constantes
    (\srcline{src/domains/impls/str_const.rs}{8}) ;
  \item \code{reference} : \code{Stability::widen} est le join
    (\src{src/domains/impls/stability.rs}{144}{146}) sur un treillis de
    hauteur finie, les ensembles \code{Versioned} étant plafonnés à
    \code{VERSIONED_LABELS_THRESHOLD = 4} labels
    (\srcline{src/domains/impls/stability.rs}{11}) ;
  \item \code{boolean}, \code{null}, \code{undef}, \code{other},
    \code{setter} : treillis finis.
\end{itemize}
Le nombre de labels est fini (fixé par l'expansion), donc le store se
stabilise, et le test $N \lleq S$ réussit. Le garde-fou à cent tours n'est
pas nécessaire à l'argument ; c'est une ceinture de sécurité.
\end{theoreme}

Les $t - 1$ premiers pas sont des joins, en nombre borné par $t$ ; le tas et
le memo store ne participent pas au test et ne peuvent pas retarder l'arrêt.
La complexité est celle de $O(k)$ tours, chacun coûtant une passe de
\code{analyze_cfg} sur le rendu, chaque effet et chaque handler, avec $k$
borné par $t$ plus la hauteur effective ci-dessus. Le dépôt ne publie pas,
au commit \snapshotCommit, de distribution de \code{iterations} sur le
corpus. Sur les exemples de ce chapitre, la valeur la plus haute est
\code{iterations = 5} (\code{Dead}), atteinte en franchissant deux seuils
avant $+\infty$.

\begin{piege}
Le garde-fou n'est jamais exercé : \code{grep 'iteration >= 100'} ne trouve
que \srcline{src/engine/fixpoint.rs}{500}, aucun test. S'il était atteint,
le résultat serait $S \widen N$ sans re-vérification : un widening est
au-dessus du join de ses opérandes, donc au-dessus de $N$, mais rien ne
garantit qu'un tour de plus depuis ce nouvel état n'ajouterait rien ; le
résultat pourrait sous-approximer d'un tour. Par le
\cref{thm:point-fixe-composant-terminaison}, ce cas ne devrait pas se
produire ; il n'est pas démontré qu'il ne se produit pas.
\end{piege}

% ===========================================================================
\section{Après la convergence}
\label{sec:point-fixe-composant-post}

Le store a convergé ; il reste à produire ce que les règles liront. Deux
passes supplémentaires et une tranche de calculs relationnels y pourvoient
(\src{src/engine/fixpoint.rs}{532}{698}).

\subsection{Rafraîchir le rendu}

\rustsnippet[label={lst:point-fixe-composant-refresh}]{src/engine/fixpoint.rs}{532}{563}{une passe de rendu de plus, sur les stores convergés}

Le commentaire dit pourquoi. Dans la boucle, le recalcul des mémos
\emph{suit} la passe de rendu : la dernière passe --- et tout ce que les
règles en lisent, \code{env_exit}, \code{block_states}, les objets alloués
dans le tas --- a donc figé les mémos de l'itération \emph{précédente}. Pour
un composant qui converge au premier tour, cela signifie $\Top$ pour chaque
\code{useCallback}, et un callback prouvablement stable se lirait \og peut
changer entre deux rendus\fg. Une passe de plus sur les stores convergés
donne aux règles ce qu'elles veulent. Deux détails : \code{inter} est
\code{None} (la passe ré-évalue et ré-alloue, elle ne ré-analyse pas les
enfants), et le store de sortie de cette passe est jeté
(\code{let (bs, _) = ...}) --- l'état convergé est $S$, pas un $S$ enrichi
d'un tour de plus.

\begin{invariant}{D'où vient chaque champ de l'\code{AnalysisResult}}{point-fixe-composant-provenance}
\code{state_store} est le $S$ convergé de la boucle. \code{block_states} et
\code{env_exit} (donc tout ce qui est évalué \og dans l'environnement de
sortie du rendu\fg{} : triggers, sites, sondes des règles) viennent de la
passe de \emph{rafraîchissement}. \code{effect_block_states} et
\code{handler_block_states} viennent de la \emph{dernière itération} de la
boucle (écrasés à chaque tour, l.~443 et~480). \code{memo_store} est le
dernier recalcul, effectué au dernier tour. \code{heap} a été muté en place
par toutes les passes, rafraîchissement et ré-exécution des effets compris.
\code{effect_setter_writes} vient de la passe décrite ci-dessous.
\end{invariant}

\subsection{Les écritures pures des effets}
\label{sec:point-fixe-composant-pure}

\rustsnippet[label={lst:point-fixe-composant-pure}]{src/engine/fixpoint.rs}{565}{594}{rejouer les effets depuis un store $\Bot$ pour isoler ce qu'ils écrivent}

Le store convergé de \code{Guarded} est $[0,10]$ ; celui de \code{Runaway}
est $[0,+\infty]$ ; tous deux sont dans \code{widen_trace}. Pour les
distinguer, il faut savoir non pas \emph{ce que vaut} le slot, mais
\emph{ce que l'effet y écrit} : $[1,10]$, borné par la garde, contre
$[1,+\infty]$. C'est ce que calcule la ré-exécution de chaque effet depuis
un store $\Bot$, dans le même \code{env_exit} et contre les mêmes mémos : le
store de sortie ne contient que les écritures des setters, sans la valeur
préexistante. Les deux portes du bras intra de \regle{infinite-loop} lisent
les deux champs dans l'ordre :

\rustsnippet[label={lst:point-fixe-composant-portes}]{src/rules/impls/infinite_loop.rs}{128}{136}{les deux portes du bras intra}

avec le prédicat de divergence :

\rustsnippet[label={lst:point-fixe-composant-unbounded}]{src/domains/impls/state_value.rs}{371}{382}{ce que \og non borné\fg{} veut dire pour une valeur du produit}

\begin{table}[htbp]\centering\small
\begin{tabularx}{\linewidth}{@{}>{\raggedright}p{0.17\linewidth}>{\raggedright}p{0.13\linewidth}>{\raggedright}p{0.1\linewidth}>{\raggedright}p{0.17\linewidth}>{\raggedright}p{0.14\linewidth}>{\raggedright\arraybackslash}X@{}}
\toprule
composant & \code{state_store} & tracé & \code{effect_setter_writes} & porte 2 & verdict \\
\midrule
\code{Runaway} & $[0,+\infty]$ & oui & $[1,+\infty]$ & non borné & \Warning \\
\code{Guarded} & $[0,10]$ & oui & $[1,10]$ & borné & absous \\
\code{Lazy} & $[7,20]$ & oui & $[8,20]$ & borné & absous \\
\code{Dead} & $[0,+\infty]$ & oui & $[1,+\infty]$ & non borné & \Warning{} (FP) \\
\code{Loader} & $[0,42]$ & non & $[42,42]$ & --- & absous, porte 1 \\
\code{MemoLoop} & $\Top$ & non & \code{ref:Versioned} & --- & absous, porte 1 (FN) \\
\code{Updater} (intra) & $[0,+\infty]$ & oui & $\Bot$ & $\Bot$ : ne pas absoudre & \Warning \\
\bottomrule
\end{tabularx}
\caption{Les deux portes de \regle{infinite-loop} sur les exemples du
chapitre (valeurs relevées par la sonde, verdicts du CLI). \code{Updater}
est donné sur le chemin intra ; le chemin du CLI est au
\cref{sec:point-fixe-composant-trous}.}
\label{tab:point-fixe-composant-portes}
\end{table}

L'\adr{14} qualifie ce bloc de \og crude precision-recovery hack\fg{} et
prévoyait de le supprimer avec le narrowing ; le narrowing ayant été
abandonné, le bloc est resté, et il est le seul mécanisme du moteur qu'un
ADR désigne ainsi. Il a un effet de bord à connaître.

\begin{piege}
Un updater fonctionnel rejoué depuis $\Bot$ écrit $\Bot$. Le premier bras de
\code{exec_setter_call} (\src{src/domains/interp/interpreter.rs}{352}{360})
lie le paramètre de \code{setX(c => c + 1)} à \code{ctx.state.get(label)}, qui
vaut $\Bot$ dans le store \code{bottom_state} ; $\Bot + 1 = \Bot$. C'est
exactement ce que montre \code{Updater} en intra :
\code{effect_setter_writes = s0=⊥} alors que \code{state_store = s0=num[0,inf]}.
La règle en tient compte : la porte 2 n'absout que si \code{!writes.is_bottom_value()}
--- $\Bot$ signifie \og je ne sais pas ce qui a été écrit\fg, pas \og rien n'a
été écrit\fg. Les tests de \file{tests/functional_updater.rs}
(\code{functional_updater_without_deps_is_flagged}, etc.) fixent ce
comportement, sur le chemin intra.
\end{piege}

\subsection{La tranche des relations}

Le reste de la fonction (\src{src/engine/fixpoint.rs}{596}{667}) ne
calcule plus rien sur le store : il dérive des \emph{relations} sur les CFG,
dans les environnements convergés. Le moteur n'est ici qu'un ordonnanceur ;
le \cref{chap:relations} détaille chaque relation. L'ordre compte :
\code{collect_hook_calls}, \code{collect_effect_info},
\code{collect_handler_info} d'abord (l.~596--598) ; puis les setters
\emph{étrangers} (un setter reçu en prop d'un parent, qualifié par son
propriétaire, l.~601--605) et les environnements par région
(\code{SiteEnvs}, l.~608--618) dont \relation{slot_writers} a besoin ;
\relation{slot_seeds} après, parce qu'elle replie les lignes de
\relation{slot_writers} et les deps des effets (l.~627--642) ;
\relation{registrations} (l.~646) ; enfin \relation{effect_triggers} :

\rustsnippet[label={lst:point-fixe-composant-triggers}]{src/engine/fixpoint.rs}{647}{667}{la relation \relation{effect_triggers}, évaluée dans \code{env_exit} contre le store convergé}

Chaque dep d'un effet est évalué par \code{eval_in_stores}
(\cref{chap:cfg-analyzer-dominance}) dans l'environnement de sortie du
rendu rafraîchi, contre \code{final_state} et le memo store, sur une copie
jetable du tas. Un dep qui \emph{est} un slot donne une ligne
\code{exact: true} --- l'effet \emph{doit} se relancer à toute écriture
fraîche, un fait must ---, un dep dont la stabilité est
\code{Versioned(S)} donne une ligne \code{exact: false} par slot de $S$ ---
l'effet \emph{peut} se relancer, un fait may (\cref{def:stabilite},
\adr{17}). Sur \code{Guarded}, \code{Fresh}, \code{Lazy} la ligne est exacte
(le dep est la variable d'état) ; sur \code{MemoLoop} elle est versionnée
(le dep est le mémo) ; sur \code{Runaway} et \code{Loader} il n'y en a pas
(pas de liste de deps, ou liste vide).

La fonction se termine en assemblant le résultat :

\rustsnippet[label={lst:point-fixe-composant-result}]{src/engine/fixpoint.rs}{670}{698}{l'\code{AnalysisResult} rendu}

\subsection{Le programme : phases 1 et 2}
\label{sec:point-fixe-composant-programme}

\code{analyze_program} (\src{src/engine/fixpoint.rs}{704}{819}) appelle
\code{analyze_component_impl} sur chaque composant, en deux phases :

\rustsnippet[label={lst:point-fixe-composant-phases}]{src/engine/fixpoint.rs}{757}{794}{la photographie de la phase 1 et le balayage intra de la phase 2}

En \emph{phase 1}, chaque racine choisie par la \code{RootStrategy}
(\code{Heuristic} : les composants qui n'apparaissent dans aucun
\code{CompApp} ; \code{--all-roots} ; \code{--entry}) est analysée avec un
\code{InterCtx}. Quand son rendu rencontre \code{<Child .../>},
\code{eval_comp_app} évalue les props, consulte un cache par égalité des
props abstraites, et sur échec appelle \code{analyze_component_inter} sur
l'enfant, dont le résultat remplace celui déjà présent dans \code{results}
(\cref{chap:transfert-stores}). Le rendu du parent étant rejoué à chaque
tour de \emph{sa} boucle, l'enfant est analysé à chaque \emph{cache miss} ;
c'est le résultat du dernier miss qui reste. Il n'y a pas de point fixe de
niveau programme : \og no separate program-level fixpoint layer\fg{}
(\adr{12}). Un enfant qui écrit un slot du parent le fait dans le
\code{SharedStateStore}, et la boucle du parent l'importe au test de
convergence suivant (\cref{lst:point-fixe-composant-test}).

En \emph{phase 2}, chaque composant que la phase 1 n'a pas atteint est
analysé \emph{intra} (\code{inter = None}, props $\Top$ par défaut
d'environnement), avec son vrai \code{ComponentId}. L'ensemble
\code{phase1_reached} est photographié \emph{avant}, pour ne pas confondre
un composant analysé en intra avec une racine (\issue{110}). Quatre
conséquences de \code{inter = None}, toutes lues dans le code et vérifiées :
\begin{enumerate}
  \item les hooks personnalisés ne sont pas greffés (l.~898) ;
  \item le \code{SummaryRegistry} n'est pas consulté non plus (le repli est
    \emph{dans} \code{expand_custom_hooks}, après ce \code{return}) : un
    \code{useQuery()} de phase 2 reste un marqueur opaque, signalé par une
    Info \regle{analysis-limit} ;
  \item tout \code{<Enfant .../>} est un enfant inconnu : \code{eval_comp_app}
    commence par \code{havoc_setter_props}
    (\src{src/domains/transfer/state_value.rs}{477}{483}), donc tout setter
    passé en prop reçoit $\Top$ --- sound ;
  \item les utilitaires sont bien greffés (cela ne dépend que de
    \code{config.function_registry}), mais la troncature de budget est
    silencieuse (\cref{sec:point-fixe-composant-expansion}).
\end{enumerate}
Le \cref{chap:inter-composants} reprend la phase 1 en détail : cache, graphe
d'appels, récursion, dépendances de rendu.

% ===========================================================================
\section{Correspondance avec React-tRace et argument de soundness}
\label{sec:point-fixe-composant-soundness}

\subsection{Un tour abstrait, quatre transitions concrètes}

L'\adr{1} fait de React-tRace la sémantique concrète du projet, avec une
justification qui concerne directement ce chapitre : \og the Tree Memory +
render loop model (StepInit $\to$ StepEffect $\to$ StepCheck) directly
corresponds to the fixpoint iteration of our abstract interpreter\fg. Le
\cref{chap:react-modele} a donné la correspondance règle par règle, dans
sa table de correspondance ; la
\cref{tab:point-fixe-composant-trace} la donne \emph{passe par passe}, avec
ce que chaque passe abstrait.

\begin{table}[htbp]\centering\small
\begin{tabularx}{\linewidth}{@{}>{\raggedright}p{0.24\linewidth}>{\raggedright}p{0.3\linewidth}>{\raggedright\arraybackslash}X@{}}
\toprule
passe de la boucle & React-tRace & ce qui est abstrait \\
\midrule
amorçage $S[\ell] \join= \sem{\mathit{init}}$ & \textsc{SttBind} (phase \textsf{Init}) & rien : une fois, hors boucle \\
passe de rendu depuis $S$ & \textsc{StepInit} puis les ré-évaluations \textsc{Succ} de \textsc{StepCheck} (\textsc{EvalOnce}/\textsc{EvalMult}) & tous les rendus possibles d'un coup ; \textsc{SttReBind} devient la lecture du store joint, convertie en \code{Versioned} (\srcline{src/domains/transfer/state_value.rs}{122}) \\
recalcul des mémos & extension esquissée au §7 de l'article & valeur = stabilité seulement \\
passes d'effets depuis \code{env_exit} & \textsc{StepEffect} : \textsc{CommitEffsPath} sur les vues à décision \textsf{Effect} & \emph{tous} les effets, ordre abstrait par le join ; fermeture = \code{env_exit} (\textsc{Eff}) ; \textsc{AppSetNormal} = mise à jour faible \\
passes de handlers & \textsc{StepEvent} (mode \emph{event loop}) & 0..N exécutions, tout ordre ; hors \code{widen_trace} \\
tranche du \code{SharedStateStore} & \textsc{AppSetNormal} vers la vue $m[p]$ d'un autre chemin & import différé au test de convergence \\
test $N \lleq S$ & \textsc{StepCheck} : \textsc{CheckNoEffect} partout $\Rightarrow$ \emph{event loop} & \og un état a changé\fg{} devient \og le store a grossi\fg \\
\code{widen_trace} & le cycle rendered $\to$ check $\to$ rendered qui ne passe jamais par \emph{event loop} (§3.1.1) & un signal abstrait : la croissance a dû être forcée à converger \\
\bottomrule
\end{tabularx}
\caption{Les passes de \code{analyze_component_impl} et les transitions de
React-tRace (\cref{chap:react-modele}).}
\label{tab:point-fixe-composant-trace}
\end{table}

Rappelons ce que le \cref{chap:react-modele} a constaté : React-tRace est la
référence \emph{conceptuelle} du moteur, pas un oracle outillé. Aucune
fonction de \file{src/engine/} ne cite une règle de l'article, et aucun test
ne compare la boucle à l'interpréteur OCaml sur les exemples du papier. La
correspondance de la \cref{tab:point-fixe-composant-trace} est donc une
lecture, faite ici, et non une propriété vérifiée par le dépôt.

\subsection{L'argument de soundness}

\begin{theoreme}{Soundness de la boucle externe (esquisse)}{point-fixe-composant-soundness}
Soit un composant, et pour tout $k$ soit $\rho_k$ l'ensemble des stores
d'état concrets (une valeur par slot) que React peut atteindre après au plus
$k$ tours --- rendu, commit, effets, puis un nombre quelconque d'événements
--- depuis le montage, pour toutes les props et toutes les suites
d'événements. Sous les hypothèses
\begin{description}
  \item[H1] les fonctions de transfert sont correctes (\cref{chap:transfert-stores}) et \code{analyze_cfg} calcule, pour chaque corps, un environnement de sortie et un store qui sur-approximent tout chemin (\cref{chap:cfg-analyzer-dominance}) ;
  \item[H2] toute écriture d'un slot passe par \code{exec_setter_call}, qui la joint au store (mise à jour faible, \src{src/domains/stores/state_store.rs}{29}{33}) ;
  \item[H3] tout corps qui peut s'exécuter est l'un des corps passés : rendu, effets, handlers, callbacks descendus par l'interpréteur ;
\end{description}
on a $\abs(\rho_k) \lleq S_k$ pour tout $k$ (en notant $S_k$ le store
après $k$ tours abstraits, prolongé par $S_{\infty}$ à l'arrêt), donc
$\abs(\bigcup_k \rho_k) \lleq S_{\infty}$ : le store convergé contient tout
état atteignable.

\emph{Esquisse.} Par récurrence sur $k$. Pour $k = 0$, $\rho_0$ est l'état
au montage, produit des initialiseurs, et $S_0$ les évalue tous
(\cref{lst:point-fixe-composant-seed}). Supposons $\abs(\rho_k) \lleq S_k$.
Un tour concret de plus part d'un état $\sigma \in \rho_k$ ; le rendu lit
$\sigma$, que $S_k$ couvre, donc par H1 l'environnement de fin de rendu
concret est dans $\gam(\mathtt{env\_exit})$ et toute écriture du rendu est
dans $R$. Les effets qui tournent (quels que soient leurs deps) partent de
cet environnement et de $\sigma$ ; chaque écriture, dans n'importe quel
ordre et quel que soit le rejeu de la file, est une valeur que
\code{analyze_cfg} sur le corps a jointe dans $E$ (H1, H2), et l'updater
lit un \code{ctx.state.get} qui couvre la file. Les handlers, exécutés
zéro ou plusieurs fois, ne produisent que des valeurs jointes dans $H$ ;
un événement supplémentaire après un tour part d'un état déjà couvert par
$N$, et $N$ est à nouveau un point de départ couvert au tour suivant, ce
qui rend compte des suites d'événements de longueur quelconque. Un setter de
ce composant appelé dans un autre composant écrit dans la tranche importée.
Donc $\abs(\rho_{k+1}) \lleq N_k \lleq S_{k+1}$, la dernière inégalité
parce que $S_{k+1}$ vaut $N_k$ ou $S_k \widen N_k \lgeq N_k$
(\code{widen_to} $\lgeq$ join, \srcline{src/domains/impls/interval.rs}{112}).
À l'arrêt, $N \lleq S$ ferme la récurrence pour tout $k$ ultérieur. $\square$
\end{theoreme}

Les trois hypothèses ne sont pas gratuites, et c'est précisément là que se
logent les trous du \cref{sec:point-fixe-composant-trous} : H2 est violée par
le second bras de \code{exec_setter_call} en inter et par le lowering de
\code{useReducer} ; H3 est violée par les fonctions de nettoyage, que la
boucle n'exécute pas. La \cref{tab:point-fixe-composant-garanties} liste,
pour mémoire, les mécanismes sur lesquels l'argument s'appuie.

\begin{table}[htbp]\centering\small
\begin{tabularx}{\linewidth}{@{}>{\raggedright}p{0.3\linewidth}>{\raggedright}p{0.27\linewidth}>{\raggedright\arraybackslash}X@{}}
\toprule
mécanisme & direction garantie & où \\
\midrule
mise à jour faible des setters & $S \lgeq$ toute valeur écrite & \src{src/domains/stores/state_store.rs}{29}{33} \\
passes depuis \code{state.clone()} & suite croissante & \srcline{src/engine/cfg_analyzer.rs}{49} \\
tous les effets à chaque tour & aucun effet oublié à cause des deps & \src{src/engine/fixpoint.rs}{418}{449} \\
handlers dans la boucle & valeurs atteignables par événements incluses & \src{src/engine/fixpoint.rs}{451}{483} \\
environnement d'effet = \code{env_exit} & l'effet voit toute valeur de fin de rendu & \srcline{src/engine/fixpoint.rs}{432} \\
\code{widen_to} $\lgeq$ join & post-point fixe & \src{src/domains/impls/interval.rs}{112}{146} \\
absent $= \Top$, join d'un seul côté $= \Top$ & pas de valeur inventée & \cref{chap:transfert-stores} \\
mémo à deps illisibles $=$ \code{Unknown} & pas de mémo figé inventé & \srcline{src/domains/transfer/state_value.rs}{67} \\
enfant inconnu : havoc des setters en props & l'enfant peut appeler tout setter reçu & \src{src/domains/transfer/state_value.rs}{477}{483} \\
budget épuisé $\to$ Info, assurances suspendues (phase 1) & pas de \code{verified} sur du code non lu & \src{src/engine/fixpoint.rs}{1620}{1628} \\
\bottomrule
\end{tabularx}
\caption{Où la soundness de la boucle est tenue.}
\label{tab:point-fixe-composant-garanties}
\end{table}

% ===========================================================================
\section{Ce que le moteur garantit, ce qu'il ne garantit pas}
\label{sec:point-fixe-composant-trous}

\file{docs/limitations.md} distingue deux registres : un \emph{défaut}
(\og the analysis returns an under-approximation or reports something
false\fg) se corrige ; un \emph{compromis} (l'analyse reste sound, elle
est seulement imprécise) se décide. Les constats qui suivent ont tous été
observés au commit \snapshotCommit{} sur les fichiers de \file{/tmp/ch13/}, et
classés selon ce critère. Plusieurs ont déjà été rencontrés en chemin ; nous
les rappelons pour que la liste soit complète.

\begin{piege}
\textbf{L'updater fonctionnel, sur le chemin du CLI (défaut, FN).}
\begin{lstlisting}[language=TypeScript,caption={\file{/tmp/ch13/updater.tsx}},label={lst:point-fixe-composant-updater}]
import { useState, useEffect } from "react";

export function Updater() {
  const [x, setX] = useState(0);
  useEffect(() => {
    setX((c) => c + 1);
  });
  return <div>{x}</div>;
}
\end{lstlisting}
Concrètement, c'est \code{Runaway} écrit avec un updater : la boucle est
infinie. En intra, le moteur la voit (\cref{tab:point-fixe-composant-portes}).
Par \code{analyze_program}, la sonde donne \code{iterations = 2},
\code{state_store = s0=⊤}, \code{widen_trace = []} et
\texttt{shared\_state = \{(ComponentId(0), 0): ref(PerRender)\}}, et le CLI :
\begin{sortie}
  [verbose] Updater: 2 iteration(s), widened: []
   1 clean component(s) hidden, rerun with --show-clean

✓  1 file(s) no issues found.
\end{sortie}
Le mécanisme est le défaut D3 du \cref{chap:transfert-stores} : le second bras
de \code{exec_setter_call} (\src{src/domains/interp/interpreter.rs}{368}{391})
s'applique aussi au propre setter du composant dès que \code{inter} est
présent, et il évalue l'argument comme une \emph{valeur} --- la fermeture
\code{c => c + 1}, une référence \code{PerRender} --- qu'il écrit dans le
\code{SharedStateStore}. La tranche est jointe à $N$
(\cref{lst:point-fixe-composant-test}), \code{c} devient \og nombre ou
référence\fg, \code{c + 1} vaut $\Top$, et le slot atteint $\Top$ en deux
tours \emph{par join} : pas de widening, \code{widen_trace} vide, porte 1
fermée. Deux causes se composent : le bras croisé ne devrait pas s'appliquer
aux slots propres (ou devrait y exécuter l'updater comme le premier bras), et
le signal de divergence ignore un slot qui atteint $\Top$ sans widening. Les
tests de \file{tests/functional_updater.rs} passent parce qu'ils appellent
\code{analyze_component}. Aucune issue ouverte ne décrit ce cas au commit
étudié.
\end{piege}

\begin{piege}
\textbf{$\Top$ atteint par join n'est pas un signal (défaut, FN).}
\code{MemoLoop} (\cref{ex:point-fixe-composant-memoloop}) boucle
concrètement et le CLI le déclare propre : le slot saute à $\Top$ au tour 0
parce que le mémo n'est pas encore calculé, et la porte 1 de
\regle{infinite-loop} ne connaît que \code{widen_trace}. Le prédicat
\code{is_unbounded}, qui compte pourtant \code{other} ($\Top$ résiduel)
comme non borné, n'est consulté qu'après la porte 1. Le bras churn lit
l'écriture comme non fraîche (\issue{157}). Deux exemples du chapitre
(\code{Updater} en inter, \code{MemoLoop}) montrent que
\og widené\fg{} et \og a divergé\fg{} ne coïncident pas : un slot peut
diverger vers $\Top$ en un ou deux joins, et ce chemin ne laisse aucune
trace.
\end{piege}

\begin{piege}
\textbf{Les fonctions de nettoyage n'écrivent pas dans le store (défaut).}
\begin{lstlisting}[language=TypeScript,caption={\file{/tmp/ch13/cleanup.tsx}},label={lst:point-fixe-composant-cleanup}]
import { useState, useEffect } from "react";

export function CleanupDeps({ id }: { id: string }) {
  const [n, setN] = useState(0);
  useEffect(() => {
    return () => setN(7);
  }, [id]);
  return <div>{n}</div>;
}
\end{lstlisting}
À chaque changement de \code{id}, React exécute le nettoyage, donc
\code{setN(7)}, puis rend avec $n = 7$. L'IR de l'effet est
\texttt{Return(FnLit \{ ... setN(7) ... \})} ; \code{exec_expr_effects} sur un
\code{FnLit} ne l'exécute pas. Sonde : \code{iterations = 0},
\code{state_store = s0=num[0,0]} --- la valeur 7 n'est pas dans le store.
C'est une sous-approximation de l'état, non répertoriée dans
\file{docs/limitations.md}. La relation \relation{slot_writers} voit la
ligne (\code{region: Effect(1), phase: Cleanup}, valeur \code{number[7, 7]},
relevé par la sonde ; \cref{chap:relations}), ce qui protège les règles qui
la lisent. Sous \file{src/rules/}, trois règles lisent \code{state_store}
directement :
\regle{redundant-set-state}, \regle{always-unstable-deps} et
\regle{wasted-subtree-render}. La première écarte d'office un slot écrit
depuis deux régions distinctes (l'ensemble \code{contested},
\src{src/rules/impls/redundant_set_state.rs}{57}{81}). Nous l'avons vérifié :
un nettoyage \code{setN(7)} plus un \code{setN(0)} dans un autre effet ou
dans un \code{onClick} ne produit aucun diagnostic. Au commit
\snapshotCommit, nous n'avons pas construit de cas où l'absence du 7 trompe
l'une des deux autres. Le trou dans le store demeure.
\end{piege}

\begin{piege}
\textbf{\code{useReducer} : l'action à la place de l'état (défaut).}
\begin{lstlisting}[language=TypeScript,caption={\file{/tmp/ch13/reducer.tsx}},label={lst:point-fixe-composant-reducer}]
import { useReducer, useEffect } from "react";

export function Reducer() {
  const [s, dispatch] = useReducer((acc: number, a: number) => acc + a, 0);
  useEffect(() => {
    dispatch(1);
  });
  return <div>{s}</div>;
}
\end{lstlisting}
Le lowering traduit \code{useReducer} en \code{HookEntry::State} et
\emph{ignore le réducteur} (\code{let _reducer = it.next(); // skip reducer fn},
\src{src/lowering/hook_extractor.rs}{715}{719}) ; \code{dispatch} est lié
comme un setter. Le moteur joint donc l'\emph{action} \code{1} au slot :
\code{s0=num[0,1]}, \code{iterations = 1}, \code{widen_trace = []}.
Concrètement $s$ vaut $0, 1, 2, 3, \dots$ et le composant boucle ; le CLI
répond \code{✓ 1 file(s) no issues found.} Pour être sound, il faudrait
exécuter le réducteur (\code{exec_body} avec \code{(état, action)}) ou
écrire $\Top$. Aucune issue ne le décrit au commit étudié.
\end{piege}

\begin{piege}
\textbf{Trois compromis déjà rencontrés.} Une branche dont la garde est
réfutée est exécutée quand même (\code{Dead},
\cref{sec:point-fixe-composant-deroules}) : faux positif, sound.
\code{WidenEvent.writers} liste tous les effets
(\cref{sec:point-fixe-composant-widen-trace}) : provenance inexacte, aucun
verdict touché. Un second appel du même hook personnalisé reste opaque, avec
un message trompeur (\cref{sec:point-fixe-composant-expansion}) : sound,
assurances suspendues. À quoi s'ajoutent deux constats sur la couche de
certification : la troncature d'inlining silencieuse en phase 2
(\cref{sec:point-fixe-composant-expansion}, un défaut au sens des
assurances publiées) et le garde-fou à cent tours jamais exercé
(\cref{sec:point-fixe-composant-terminaison}).
\end{piege}

\begin{table}[htbp]\centering\small
\begin{tabularx}{\linewidth}{@{}>{\raggedright}p{0.36\linewidth}>{\raggedright}p{0.17\linewidth}>{\raggedright}p{0.15\linewidth}>{\raggedright\arraybackslash}X@{}}
\toprule
constat & exemple & registre & référence \\
\midrule
updater fonctionnel écrit comme fermeture dans le store partagé & \code{Updater} & défaut (FN) & \cref{chap:transfert-stores} D3 ; aucune issue \\
$\Top$ par join sans widening : porte 1 fermée & \code{MemoLoop}, \code{Updater} & défaut (FN) & \issue{157}, \issue{144} \\
nettoyages hors du store & \code{CleanupDeps} & défaut & non documenté \\
\code{useReducer} : action jointe, réducteur ignoré & \code{Reducer} & défaut & non documenté \\
troncature d'inlining silencieuse en phase 2 & \code{A}/\code{B} & défaut (assurances) & non documenté \\
branche réfutée exécutée & \code{Dead} & compromis (FP) & \cref{chap:cfg-analyzer-dominance} \\
\code{writers} = tous les effets & \code{Mutual} & compromis (provenance) & \cref{sec:point-fixe-composant-widen-trace} \\
second appel d'un hook personnalisé opaque & \code{TwoClocks} & compromis & \srcline{src/engine/fixpoint.rs}{884} \\
\code{useState(makeInit)} amorcé à $\Top$ & \code{LazyVar} & compromis (FP) & \cref{sec:point-fixe-composant-amorcage} \\
cap à 100 sans re-vérification ni test & --- & risque théorique & \srcline{src/engine/fixpoint.rs}{500} \\
\bottomrule
\end{tabularx}
\caption{Les constats du chapitre, classés selon les deux registres de
\file{docs/limitations.md}.}
\label{tab:point-fixe-composant-constats}
\end{table}

% ===========================================================================
\begin{resume}
\begin{itemize}
  \item \code{analyze_component_impl} (\file{src/engine/fixpoint.rs}) ne
    simule pas React : elle calcule un post-point fixe du \code{StateStore},
    un sur-ensemble des états atteignables par tous les tours et toutes les
    suites d'événements. Un tour abstrait représente tous les tours concrets
    de même rang.
  \item \emph{Préparation}, une fois : constantes de module dans
    \code{initial_env} et \code{module_env} ; greffes des utilitaires puis
    des hooks personnalisés (\code{SpliceIds} monotone, garde par nom,
    résumés par \code{retag_marker}) ; seuils = tous les littéraux
    numériques ; \code{callback_bodies} ; amorçage des \code{useState},
    initialiseur paresseux exécuté.
  \item \emph{Un tour} : rendu depuis \code{initial_env} ; mémos recalculés
    depuis \code{env_exit} (valeur = stabilité, $\Top$ au tour 0, hors test
    de convergence) ; tous les effets depuis \code{env_exit} et le même
    store (deps ignorées, batching abstrait) ; handlers de même, hors
    signal ; $N = R \join E \join H \join$ tranche partagée ; arrêt si
    $N \lleq S$, ce qui équivaut à $N = S$.
  \item \emph{Convergence} : compteur global des tours qui changent l'état ;
    joins jusqu'à \code{widen_threshold = 3}, puis \code{widen_to} à seuils
    sur tout le store, \code{widen_trace} sur les seuls labels changés par
    $R \join E$ (première itération retenue) ; garde-fou à 100. Terminaison
    par hauteur finie de chaque composante de \code{StateValue}.
  \item \emph{Post-convergence} : passe de rafraîchissement (mémos à jour
    pour les règles, \code{inter = None}) ; \code{effect_setter_writes}
    depuis $\Bot$ (porte 2 de \regle{infinite-loop} ; un updater y écrit
    $\Bot$) ; tranche des relations, dont \relation{effect_triggers}
    (exact = must, versionné = may).
  \item \emph{Deux moteurs} : intra (\code{analyze_component}, phase 2) sans
    greffe de hooks ni résumés ni store partagé ; inter (phase 1) avec.
  \item \emph{Trous au commit \snapshotCommit} : updater en inter (FN),
    $\Top$ par join (FN), nettoyages et \code{useReducer} hors du store,
    troncature silencieuse en phase 2 ; compromis : branche réfutée
    exécutée, \code{writers} imprécis, second hook opaque.
\end{itemize}
\begin{sloppypar}
Types et fichiers rencontrés : \code{Config} et \code{SpliceIds}
(\file{src/engine/fixpoint.rs}) ; \code{AnalysisResult} et \code{WidenEvent}
(\file{src/engine/analysis_result.rs}) ; \code{StateStore}, \code{MemoStore}
et \code{SharedStateStore} (\file{src/domains/stores/}) ;
\code{EffectTrigger} (\file{src/engine/triggers.rs}) ; \code{HookRegistry}
et \code{FunctionRegistry}. Le \cref{chap:relations} reprend la dernière
tranche de la fonction : ce que sont \relation{slot_writers},
\relation{slot_seeds}, \relation{registrations} et \relation{effect_triggers},
et ce qu'une règle a le droit d'en certifier.
\end{sloppypar}
\end{resume}

\section*{Exercices}
\addcontentsline{toc}{section}{Exercices}

\begin{exercice}{Un compteur gardé à pas de deux}{point-fixe-composant-exo-pas2}
Dérouler la boucle externe pour
\texttt{useEffect(() => \{ if (count < 3) setCount(count + 2); \}, [count])}
avec \code{useState(0)}. Quels sont les seuils ? Quelle valeur finale, quel
\code{iterations}, quel \code{widen_trace} ? Vérifier avec le CLI.
\end{exercice}
\begin{solution}
Seuils $\{0, 2, 3\}$. Tour 0 : dans la branche, \code{count} est raffiné à
$[0,0]$, l'effet écrit $[2,2]$, $N = [0,2]$. Tour 1 : \code{count} $= [0,2]$,
raffiné par $< 3$ à $[0,2]$, écriture $[2,4]$, $N = [0,4]$. Tour 2 :
\code{count} $= [0,4]$, raffiné à $[0,2]$, écriture $[2,4]$, $N = [0,4] \lleq S$.
Arrêt sans widening : \code{iterations = 2}, \code{s0=num[0,4]},
\code{widen_trace = []}, \code{effect_setter_writes = s0=num[2,4]} (sonde,
\file{/tmp/ch13/exo1.tsx}). CLI : \code{Exo: 2 iteration(s), widened: []},
aucun diagnostic. Le raffinement de garde a suffi ; le seuil 3 n'a pas
servi.
\end{solution}

\begin{exercice}{Seuil bas}{point-fixe-composant-exo-seuil}
Refaire \code{Counter} (\cref{ex:point-fixe-composant-counter}) et
\code{Runaway} avec \code{widen_threshold = 1}. Que devient chaque tableau ?
Comparer au test \code{widened_labels_triggered_with_low_threshold}
(\srcline{src/engine/fixpoint.rs}{2011}).
\end{exercice}

\begin{exercice}{Égalité}{point-fixe-composant-exo-egalite}
La \cref{prop:point-fixe-composant-egalite} conclut que le test
\code{new_state.leq(&state)} équivaut à $N = S$. Où l'implémentation
\code{StateStore::leq} (\cref{lst:point-fixe-composant-leq}) diffère-t-elle
d'un test d'égalité structurelle, et pourquoi la différence est-elle sans
effet dans la boucle ? (Considérer un label présent dans $S$ et absent de
$N$, et deux valeurs incomparables.)
\end{exercice}

\begin{exercice}{Provenance}{point-fixe-composant-exo-writers}
Construire un composant où deux effets écrivent deux slots distincts,
observer \code{WidenEvent.writers} avec \code{--trace}, puis écrire le
correctif de la l.~444 de \file{src/engine/fixpoint.rs} qui ne retient que
les labels réellement changés par la passe. Quel test du dépôt faudrait-il
ajouter ?
\end{exercice}

\begin{exercice}{Zéro itération}{point-fixe-composant-exo-fresh}
Expliquer pourquoi \code{Fresh} (\cref{lst:point-fixe-composant-fresh})
converge avec \code{iterations = 0}, ce que valent \code{widen_trace} et
\code{effect_setter_writes}, et par quelle chaîne de faits l'\Error{} est
tout de même produite (lire les \cref{chap:relations,chap:churn}).
\end{exercice}

\begin{exercice}{Intra contre inter}{point-fixe-composant-exo-updater}
Reproduire \code{Updater} (\cref{lst:point-fixe-composant-updater}) avec
\code{analyze_component} puis avec \code{analyze_program}, et suivre la
valeur écrite dans le \code{SharedStateStore} depuis
\src{src/domains/interp/interpreter.rs}{368}{391} jusqu'à
\src{src/engine/fixpoint.rs}{488}{493}. Proposer un correctif au niveau
central, et dire lequel des deux défauts composés il corrige.
\end{exercice}

\begin{exercice}{Terminaison}{point-fixe-composant-exo-terminaison}
Écrire l'argument de terminaison complet en donnant, pour chaque composante
de \code{StateValue}, une borne explicite sur la longueur d'une chaîne
$v \lleq v \widen_T v' \lleq \dots$ en fonction de $|T|$,
\code{STR_WIDEN_THRESHOLD} et \code{VERSIONED_LABELS_THRESHOLD}. En déduire
une borne sur \code{iterations}, et la comparer au cap de 100.
\end{exercice}
